\documentclass[aspectratio=169]{beamer}\usepackage[]{graphicx}\usepackage[]{xcolor}
% maxwidth is the original width if it is less than linewidth
% otherwise use linewidth (to make sure the graphics do not exceed the margin)
\makeatletter
\def\maxwidth{ %
  \ifdim\Gin@nat@width>\linewidth
    \linewidth
  \else
    \Gin@nat@width
  \fi
}
\makeatother

\definecolor{fgcolor}{rgb}{0.345, 0.345, 0.345}
\newcommand{\hlnum}[1]{\textcolor[rgb]{0.686,0.059,0.569}{#1}}%
\newcommand{\hlsng}[1]{\textcolor[rgb]{0.192,0.494,0.8}{#1}}%
\newcommand{\hlcom}[1]{\textcolor[rgb]{0.678,0.584,0.686}{\textit{#1}}}%
\newcommand{\hlopt}[1]{\textcolor[rgb]{0,0,0}{#1}}%
\newcommand{\hldef}[1]{\textcolor[rgb]{0.345,0.345,0.345}{#1}}%
\newcommand{\hlkwa}[1]{\textcolor[rgb]{0.161,0.373,0.58}{\textbf{#1}}}%
\newcommand{\hlkwb}[1]{\textcolor[rgb]{0.69,0.353,0.396}{#1}}%
\newcommand{\hlkwc}[1]{\textcolor[rgb]{0.333,0.667,0.333}{#1}}%
\newcommand{\hlkwd}[1]{\textcolor[rgb]{0.737,0.353,0.396}{\textbf{#1}}}%
\let\hlipl\hlkwb

\usepackage{framed}
\makeatletter
\newenvironment{kframe}{%
 \def\at@end@of@kframe{}%
 \ifinner\ifhmode%
  \def\at@end@of@kframe{\end{minipage}}%
  \begin{minipage}{\columnwidth}%
 \fi\fi%
 \def\FrameCommand##1{\hskip\@totalleftmargin \hskip-\fboxsep
 \colorbox{shadecolor}{##1}\hskip-\fboxsep
     % There is no \\@totalrightmargin, so:
     \hskip-\linewidth \hskip-\@totalleftmargin \hskip\columnwidth}%
 \MakeFramed {\advance\hsize-\width
   \@totalleftmargin\z@ \linewidth\hsize
   \@setminipage}}%
 {\par\unskip\endMakeFramed%
 \at@end@of@kframe}
\makeatother

\definecolor{shadecolor}{rgb}{.97, .97, .97}
\definecolor{messagecolor}{rgb}{0, 0, 0}
\definecolor{warningcolor}{rgb}{1, 0, 1}
\definecolor{errorcolor}{rgb}{1, 0, 0}
\newenvironment{knitrout}{}{} % an empty environment to be redefined in TeX

\usepackage{alltt}

% Set lecture number for later use


% Part common to all the lectures
\subtitle{MATH 2740 -- Mathematics of Data Science -- Unit U11}
\author{\texorpdfstring{Julien Arino\newline\url{julien.arino@umanitoba.ca}}{Julien Arino}}
\institute{Department of Mathematics @ University of Manitoba}

% Title of the lecture
\title{Vector space models, sparse text matrices and NMF}



\usetheme{default}
% Slide setup, colour independent

\usepackage{amsmath,amssymb,amsthm}
\usepackage[utf8]{inputenc}
\usepackage{colortbl}
\usepackage{bm}
\usepackage{xcolor}
\usepackage{dsfont}
\usepackage{setspace}
% To use \ding{234} and the like
\usepackage{pifont}
% To cross reference between slide files
\usepackage{zref-xr,zref-user}
% Use something like
% \zexternaldocument{fileI}
% in the tex files. And cite using \zref instead of \ref

% Cross-reference system - see CROSS-REFERENCE-SETUP.md for manual setup instructions
\usepackage{booktabs}
\usepackage{marvosym}
\usepackage{cancel}
%\usepackage{transparent}
% Make doi clickable in the bibliography?
\usepackage{doi}

\usepackage[T1]{fontenc}

\usepackage{longtable}

% For heavier titles
\usepackage{helvet} % Enables Helvetica font family


% Fields and the like
\def\IC{\mathbb{C}}
\def\IE{\mathbb{E}}
\def\IF{\mathbb{F}}
\def\II{\mathbb{I}}
\def\IJ{\mathbb{J}}
\def\IK{\mathbb{K}}
\def\IM{\mathbb{M}}
\def\IN{\mathbb{N}}
\def\IP{\mathbb{P}}
\def\IR{\mathbb{R}}
\newcommand{\IRplus}{\mathbb{R}_{\ge 0}}
\def\IZ{\mathbb{Z}}
\def\11{\mathds{1}}


% Bold lowercase
\def\ba{\bm{a}}
\def\bb{\bm{b}}
\def\bc{\bm{c}}
\def\bd{\bm{d}}
\def\be{\bm{e}}
\def\bf{\bm{f}}
\def\bg{\bm{g}}
\def\bh{\bm{h}}
\def\bi{\bm{i}}
\def\bj{\bm{j}}
\def\bk{\bm{k}}
\def\bn{\bm{n}}
\def\bp{\bm{p}}
\def\br{\bm{r}}
\def\bs{\bm{s}}
\def\bu{\bm{u}}
\def\bv{\bm{v}}
\def\bw{\bm{w}}
\def\bx{\bm{x}}
\def\by{\bm{y}}
\def\bz{\bm{z}}
\newcommand{\vect}[1]{\bm{#1}}

% Bold capitals
\def\bB{\bm{B}}
\def\bD{\bm{D}}
\def\bE{\bm{E}}
\def\bF{\bm{F}}
\def\bG{\bm{G}}
\def\bI{\bm{I}}
\def\bL{\bm{L}}
\def\bN{\bm{N}}
\def\bP{\bm{P}}
\def\bR{\bm{R}}
\def\bS{\bm{S}}
\def\bT{\bm{T}}
\def\bX{\bm{X}}

% Bold numbers
\def\b0{\bm{0}}

% Bold greek
\bmdefine{\bmu}{\bm{\mu}}
\def\bphi{\bm{\phi}}
\def\bvarphi{\bm{\varphi}}
\def\bPi{\bm{\Pi}}
\def\bGamma{\bm{\Gamma}}

% Bold red sentence
\def\boldred#1{{\color{red}\textbf{#1}}}
\def\defword#1{{\color{orange}\textbf{#1}}}

% Caligraphic letters
\def\A{\mathcal{A}}
\def\B{\mathcal{B}}
\def\C{\mathcal{C}}
\def\D{\mathcal{D}}
\def\E{\mathcal{E}}
\def\F{\mathcal{F}}
\def\G{\mathcal{G}}
\def\H{\mathcal{H}}
\def\I{\mathcal{I}}
\def\L{\mathcal{L}}
\def\M{\mathcal{M}}
\def\N{\mathcal{N}}
\def\P{\mathcal{P}}
\def\R{\mathcal{R}}
\def\S{\mathcal{S}}
\def\T{\mathcal{T}}
\def\U{\mathcal{U}}
\def\V{\mathcal{V}}

% Adding space for prime (') where needed
\def\pprime{\,'}
% Adding space for star (\star) where needed
\def\pstar{{\,\star}}

% tt font for code
\def\code#1{{\tt #1}}

% i.e., e.g.
\def\eg{\emph{e.g.}}
\def\ie{\emph{i.e.}}


% Operators and special symbols
\def\nbOne{{\mathchoice {\rm 1\mskip-4mu l} {\rm 1\mskip-4mu l}
{\rm 1\mskip-4.5mu l} {\rm 1\mskip-5mu l}}}
\def\cov{\ensuremath{\mathsf{cov}}}
\def\Var{\ensuremath{\mathsf{Var}\ }}
\def\Im{\textrm{Im}\;}
\def\Re{\textrm{Re}\;}
\def\det{\ensuremath{\mathsf{det}}}
\def\diag{\ensuremath{\mathsf{diag}}}
\def\nullspace{\ensuremath{\mathsf{null}}}
\def\nullity{\ensuremath{\mathsf{nullity}}}
\def\rank{\ensuremath{\mathsf{rank}}}
\def\range{\ensuremath{\mathsf{range}}}
\def\sgn{\ensuremath{\mathsf{sgn}}}
\def\Span{\ensuremath{\mathsf{span}}}
\def\tr{\ensuremath{\mathsf{tr}}}
\def\imply{$\Rightarrow$}
\def\restrictTo#1#2{\left.#1\right|_{#2}}
\newcommand{\parallelsum}{\mathbin{\!/\mkern-5mu/\!}}
\def\dsum{\mathop{\displaystyle \sum }}%
\def\dind#1#2{_{\substack{#1\\ #2}}}

\newcommand{\Qmatrix}[1]{%
  \begin{pmatrix}#1\end{pmatrix}%
}

\DeclareMathOperator{\GL}{GL}
\DeclareMathOperator{\Rel}{Re}
\def\Nt#1{\left|\!\left|\!\left|#1\right|\!\right|\!\right|}
\newcommand{\tripbar}{|\! |\! |}



% The beamer bullet (in base colour)
\def\bbullet{\leavevmode\usebeamertemplate{itemize item}\ }

% Theorems and the like
\newtheorem{importantdefinition}[theorem]{Definition}
\newtheorem{proposition}[theorem]{Proposition}
\newtheorem{property}[theorem]{Property}
\newtheorem{importantproperty}[theorem]{Property}
\newtheorem{importanttheorem}[theorem]{Theorem}
%\newtheorem{lemma}[theorem]{Lemma}
%\newtheorem{corollary}[theorem]{Corollary}
\newtheorem{remark}[theorem]{Remark}
\setbeamertemplate{theorems}[numbered]
%\setbeamertemplate{theorems}[ams style]

%
%\usecolortheme{orchid}
%\usecolortheme{orchid}

\def\red{\color[rgb]{1,0,0}}
\def\blue{\color[rgb]{0,0,1}}
\def\green{\color[rgb]{0,1,0}}

% Fix skipping lines after items in the bibliography
\setbeamertemplate{bibliography entry title}{}
\setbeamertemplate{bibliography entry location}{}
\setbeamertemplate{bibliography entry note}{}

% Get rid of navigation stuff
\setbeamertemplate{navigation symbols}{}

% Set footline/header line
\setbeamertemplate{footline}
{%
\quad p. \insertpagenumber \quad--\quad \insertsection\vskip2pt
}
% \setbeamertemplate{headline}
% {%
% \quad\insertsection\hfill p. \insertpagenumber\quad\mbox{}\vskip2pt
% }


\makeatletter
\newlength\beamerleftmargin
\setlength\beamerleftmargin{\Gm@lmargin}
\makeatother

% Colours for special pages
\def\extraContent{yellow!20}


%%%%%%%%%%%%%%%%%
\usepackage{tikz}
\usetikzlibrary{shapes,arrows}
\usetikzlibrary{positioning}
\usetikzlibrary{shapes.symbols,shapes.callouts,patterns}
\usetikzlibrary{calc,fit}
\usetikzlibrary{backgrounds}
\usetikzlibrary{decorations.pathmorphing,fit,petri}
\usetikzlibrary{automata}
\usetikzlibrary{fadings}
\usetikzlibrary{patterns,hobby}
\usetikzlibrary{backgrounds,fit,petri}
\usetikzlibrary{tikzmark}

\usepackage{pgfplots}
\pgfplotsset{compat=1.6}
\pgfplotsset{ticks=none}

\usetikzlibrary{decorations.markings}
\usetikzlibrary{arrows.meta}
\tikzset{>=stealth}

% For tikz
\tikzstyle{cloud} = [draw, ellipse,fill=red!20, node distance=0.87cm,
minimum height=2em]
\tikzstyle{line} = [draw, -latex']


%%% For max frame images
\newenvironment{changemargin}[2]{%
\begin{list}{}{%
\setlength{\topsep}{0pt}%
\setlength{\leftmargin}{#1}%
\setlength{\rightmargin}{#2}%
\setlength{\listparindent}{\parindent}%
\setlength{\itemindent}{\parindent}%
\setlength{\parsep}{\parskip}%
}%
\item[]}{\end{list}}


% Make one image take up the entire slide content area in beamer,.:
% centered/centred full-screen image, with title:
% This uses the whole screen except for the 1cm border around it
% all. 128x96mm
\newcommand{\titledFrameImage}[2]{
\begin{frame}{#1}
%\begin{changemargin}{-1cm}{-1cm}
\begin{center}
\includegraphics[width=108mm,height=\textheight,keepaspectratio]{#2}
\end{center}
%\end{changemargin}
\end{frame}
}

% Make one image take up the entire slide content area in beamer.:
% centered/centred full-screen image, no title:
% This uses the whole screen except for the 1cm border around it
% all. 128x96mm
\newcommand{\plainFrameImage}[1]{
\begin{frame}[plain]
%\begin{changemargin}{-1cm}{-1cm}
\begin{center}
\includegraphics[width=108mm,height=76mm,keepaspectratio]{#1}
\end{center}
%\end{changemargin}
\end{frame}
}

% Make one image take up the entire slide area, including borders, in beamer.:
% centered/centred full-screen image, no title:
% This uses the entire whole screen
\newcommand{\maxFrameImage}[1]{
\begin{frame}[plain]
\begin{changemargin}{-1cm}{-1cm}
\begin{center}
\includegraphics[width=\paperwidth,height=\paperheight,keepaspectratio]
{#1}
\end{center}
\end{changemargin}
\end{frame}
}

% This uses the entire whole screen (to include in frame)
\newcommand{\maxFrameImageNoFrame}[1]{
\begin{changemargin}{-1cm}{-1cm}
\begin{center}
\includegraphics[width=\paperwidth,height=0.99\paperheight,keepaspectratio]
{#1}
\end{center}
\end{changemargin}
}

% Make one image take up the entire slide area, including borders, in beamer.:
% centered/centred full-screen image, no title:
% This uses the entire whole screen
\newcommand{\maxFrameImageColor}[2]{
\begin{frame}[plain]
\setbeamercolor{normal text}{bg=#2!20}
\begin{changemargin}{-1cm}{-1cm}
\begin{center}
\includegraphics[width=\paperwidth,height=\paperheight,keepaspectratio]
{#1}
\end{center}
\end{changemargin}
\end{frame}
}


\usepackage{tikz}
\usetikzlibrary{patterns,hobby,matrix}
\usepackage{pgfplots}
\pgfplotsset{compat=1.6}
\pgfplotsset{ticks=none}

\usetikzlibrary{backgrounds}
\usetikzlibrary{decorations.markings}
\usetikzlibrary{arrows.meta}
\tikzset{>=stealth}

\tikzset{
  clockwise arrows/.style={
    postaction={
      decorate,
      decoration={
        markings,
        mark=between positions 0.1 and 0.9 step 40pt with {\arrow{>}},
   }}}}


% New integrated section command: creates section and section slide
\newcommand{\Ssection}[2]{
\section{#1}
\begin{frame}[noframenumbering,plain]
  \begin{tikzpicture}[remember picture,overlay]
    \node[above right,inner sep=0pt,opacity=0.2] at (current page.south west)
    {
        \includegraphics[height=\paperheight,width=\paperwidth]{#2}
    };
  \end{tikzpicture}
  \setbeamercolor{section in toc}{fg=section_page_list_colour}
  \setbeamerfont{section in toc}{size=\Large,series=\bfseries}
  \setbeamertemplate{section in toc shaded}[default][60]
  \tableofcontents[
    currentsection,
    sectionstyle=show/shaded,
    subsectionstyle=show/hide/hide,
    subsubsectionstyle=hide/hide/hide]
\end{frame}
\addtocounter{page}{-1}
}

% New integrated section command with subsections: creates section and section slide showing subsections
\newcommand{\SsectionWithSubs}[2]{
\section{#1}
\begin{frame}[noframenumbering,plain]
  \begin{tikzpicture}[remember picture,overlay]
    \node[above right,inner sep=0pt,opacity=0.2] at (current page.south west)
    {
        \includegraphics[height=\paperheight,width=\paperwidth]{#2}
    };
  \end{tikzpicture}
  \setbeamercolor{section in toc}{fg=section_page_list_colour}
  \setbeamerfont{section in toc}{size=\Large,series=\bfseries}
  \setbeamertemplate{section in toc shaded}[default][60]
  \tableofcontents[
    currentsection,
    sectionstyle=show/hide,
    subsectionstyle=show/show/hide,
    subsubsectionstyle=hide/hide/hide]
\end{frame}
\addtocounter{page}{-1}
}

% New integrated subsection command: creates subsection and subsection slide
\newcommand{\Ssubsection}[2]{
\subsection{#1}
\begin{frame}[noframenumbering,plain]
  \begin{tikzpicture}[remember picture,overlay]
    \node[above right,inner sep=0pt,opacity=0.2] at (current page.south west)
    {
        \includegraphics[height=\paperheight,width=\paperwidth]{#2}
    };
  \end{tikzpicture}
  \setbeamercolor{section in toc}{fg=subsection_page_list_colour}
  \setbeamerfont{section in toc}{size=\Large,series=\bfseries}
  \setbeamertemplate{section in toc shaded}[default][60]
  \setbeamerfont{subsection in toc}{series=\bfseries}
  \setbeamertemplate{subsection in toc shaded}[default][50]
  \tableofcontents[
    currentsection,
    sectionstyle=show/hide,
    subsectionstyle=show/shaded/hide,
    subsubsectionstyle=hide/hide/hide]
\end{frame}
\addtocounter{page}{-1}
}

% New integrated subsubsection command: creates subsubsection and subsubsection slide
\newcommand{\Ssubsubsection}[2]{
\subsubsection{#1}
\begin{frame}[noframenumbering,plain]
  \begin{tikzpicture}[remember picture,overlay]
    \node[above right,inner sep=0pt,opacity=0.2] at (current page.south west)
    {
        \includegraphics[height=\paperheight,width=\paperwidth]{#2}
    };
  \end{tikzpicture}
  \setbeamercolor{section in toc}{fg=subsub_header_section}
  \setbeamerfont{section in toc}{size=\Large,series=\bfseries}
  \setbeamertemplate{section in toc shaded}[default][60]
  \setbeamerfont{subsection in toc}{series=\bfseries}
  \setbeamertemplate{subsection in toc shaded}[default][50]
  \setbeamertemplate{subsubsection in toc shaded}[default][50]
  \tableofcontents[
    currentsection,
    sectionstyle=show/hide,
    subsectionstyle=show/hide/hide,
    subsubsectionstyle=show/shaded/hide]
\end{frame}
\addtocounter{page}{-1}
}

% Legacy commands (kept for backward compatibility)
% Beginning of a section
\newcommand{\newSectionSlide}[1]{
\begin{frame}[noframenumbering,plain]
  \begin{tikzpicture}[remember picture,overlay]
    \node[above right,inner sep=0pt,opacity=0.2] at (current page.south west)
    {
        \includegraphics[height=\paperheight,width=\paperwidth]{#1}
    };
  \end{tikzpicture}
  \setbeamercolor{section in toc}{fg=section_page_list_colour}
  \setbeamerfont{section in toc}{size=\Large,series=\bfseries}
  \setbeamertemplate{section in toc shaded}[default][60]
  \tableofcontents[
    currentsection,
    sectionstyle=show/shaded,
    subsectionstyle=show/hide/hide,
    subsubsectionstyle=hide/hide/hide]
\end{frame}
\addtocounter{page}{-1}
}

% Beginning of a section in which we also show subsections
\newcommand{\newSectionWithSubsSlide}[1]{
	\begin{frame}[noframenumbering,plain]
		\begin{tikzpicture}[remember picture,overlay]
			\node[above right,inner sep=0pt,opacity=0.2] at (current page.south west)
			{
				\includegraphics[height=\paperheight,width=\paperwidth]{#1}
			};
		\end{tikzpicture}
		\setbeamercolor{section in toc}{fg=section_page_list_colour}
		\setbeamerfont{section in toc}{size=\Large,series=\bfseries}
		\setbeamertemplate{section in toc shaded}[default][60]
		\tableofcontents[
		currentsection,
		sectionstyle=show/hide,
		subsectionstyle=show/show/hide,
		subsubsectionstyle=hide/hide/hide]
	\end{frame}
	\addtocounter{page}{-1}
}

% Beginning of a subsection
\newcommand{\newSubSectionSlide}[1]{
\begin{frame}[noframenumbering,plain]
  \begin{tikzpicture}[remember picture,overlay]
    \node[above right,inner sep=0pt,opacity=0.2] at (current page.south west)
    {
        \includegraphics[height=\paperheight,width=\paperwidth]{#1}
    };
  \end{tikzpicture}
  \setbeamercolor{section in toc}{fg=subsection_page_list_colour}
  \setbeamerfont{section in toc}{size=\Large,series=\bfseries}
  \setbeamertemplate{section in toc shaded}[default][60]
  \setbeamerfont{subsection in toc}{series=\bfseries}
  \setbeamertemplate{subsection in toc shaded}[default][50]
  \tableofcontents[
    currentsection,
    sectionstyle=show/hide,
    subsectionstyle=show/shaded/hide,
    subsubsectionstyle=hide/hide/hide]
\end{frame}
\addtocounter{page}{-1}
}

% Beginning of a subsubsection
\newcommand{\newSubSubSectionSlide}[1]{
\begin{frame}[noframenumbering,plain]
  \begin{tikzpicture}[remember picture,overlay]
    \node[above right,inner sep=0pt,opacity=0.2] at (current page.south west)
    {
        \includegraphics[height=\paperheight,width=\paperwidth]{#1}
    };
  \end{tikzpicture}
  \setbeamercolor{section in toc}{fg=subsub_header_section}
  \setbeamerfont{section in toc}{size=\Large,series=\bfseries}
  \setbeamertemplate{section in toc shaded}[default][60]
  \setbeamerfont{subsection in toc}{series=\bfseries}
  \setbeamertemplate{subsection in toc shaded}[default][50]
  \setbeamertemplate{subsubsection in toc shaded}[default][50]
  \tableofcontents[
    currentsection,
    sectionstyle=show/hide,
    subsectionstyle=show/hide/hide,
    subsubsectionstyle=show/shaded/hide]
\end{frame}
\addtocounter{page}{-1}
}


   %%%%%%%%%%%
% To have links to parts in the outline
\makeatletter
\AtBeginPart{%
  \addtocontents{toc}{\protect\beamer@partintoc{\the\c@part}{\beamer@partnameshort}{\the\c@page}}%
}
%% number, shortname, page.
\providecommand\beamer@partintoc[3]{%
  \ifnum\c@tocdepth=-1\relax
    % requesting onlyparts.
    \makebox[6em]{Part #1:} \textcolor{green!30!blue}{\hyperlink{#2}{#2}}
    \par
  \fi
}
\define@key{beamertoc}{onlyparts}[]{%
  \c@tocdepth=-1\relax
}
\makeatother%

\newcommand{\nameofthepart}{}
\newcommand{\nupart}[1]%
    {   \part{#1}%
        \renewcommand{\nameofthepart}{#1}%
        {
          \setbeamercolor{background canvas}{bg=orange!50}
          \begin{frame}{#1}%\partpage 
          \hypertarget{\nameofthepart}{}\tableofcontents%
          \end{frame}
        }
    }

% This command creates a title page using TikZ only
\newcommand{\tikztitlepage}[1]{%
\begin{frame}[plain,noframenumbering]
  \begin{tikzpicture}[remember picture,overlay]
    % Background image
    \node[above right,inner sep=0pt,opacity=0.1] 
      at (current page.south west) 
      {\includegraphics[width=\paperwidth,height=\paperheight]{#1}};

    % University logo
    \node[anchor=north east, inner sep=5pt, opacity=0.9] 
      at (current page.north east)
      {\includegraphics[width=0.2\textwidth]{FIGS-slides-admin/UM-logo-horizontal-CMYK.png}};
    
    % Title
    \node[anchor=center, align=center, 
          font=\fontsize{13}{15}\bfseries\color{UMbrown}, 
          text width=0.9\textwidth] 
          at ([yshift=2cm]current page.center)
          {\inserttitle};

      % Authors
      \node[anchor=center, align=center,
        font=\fontsize{10}{12}\bfseries\color{UMbrown},
        text width=0.7\textwidth]
        at ([yshift=0.8cm]current page.center)
        {\insertauthor};

      % Affiliation
      \node[anchor=north, align=center,
        font=\fontsize{9}{11}\color{UMbrown},
        text width=0.7\textwidth]
        at ([yshift=-0.2cm]current page.center)
        {\insertaffiliation};      
    % Date
    \node[anchor=north, align=center, 
          font=\fontsize{12}{16}\bfseries\color{UMbrown},
          text width=0.7\textwidth] 
          at ([yshift=0.2cm]current page.center)
          {\insertdate};

    % Land acknowledgement
    \node[anchor=south, align=justify, 
          font=\footnotesize, text=black, 
          text width=1.1\textwidth] 
          at ([yshift=0.5cm]current page.south)
          {The University of Manitoba campuses are located on original lands of Anishinaabeg, Ininew, Anisininew, Dakota and Dene peoples, and on the National Homeland of the Red River Métis.\\
          We respect the Treaties that were made on these territories, we acknowledge the harms and mistakes of the past, and we dedicate ourselves to move forward in partnership with Indigenous communities in a spirit of Reconciliation and collaboration.};
  \end{tikzpicture}
  \addtocounter{page}{-1}
\end{frame}
}
% The title page with figure
% \newcommand{\titlepagewithfigure}[1]{%
%   \begin{frame}[noframenumbering,plain]
%     \begin{tikzpicture}[remember picture,overlay]
%       \node[above right,inner sep=0pt,opacity=0.1] at (current page.south west)
%       {
%           \includegraphics[height=\paperheight,width=\paperwidth]{#1}
%       };
%       \node[anchor=north east,
%       inner sep=5pt,
%       opacity=0.9] at (current page.north east)
%       {
%           \includegraphics[width=0.2\textwidth]{FIGS-slides-admin/UM-logo-horizontal-CMYK.png}
%       };
%       \node[anchor=south, 
%       align=justify, 
%       text=black, 
%       text width=1.1\textwidth,
%       font=\footnotesize]  (land_acknowledgement)
%       at (current page.south) 
%       {The University of Manitoba campuses are located on original lands of Anishinaabeg, Ininew, Anisininew, Dakota and Dene peoples, and on the National Homeland of the Red River Métis.
%       We respect the Treaties that were made on these territories, we acknowledge the harms and mistakes of the past, and we dedicate ourselves to move forward in partnership with Indigenous communities in a spirit of Reconciliation and collaboration.};  
%       % \node[align=center, anchor=south,
%       % above=0.5cm of land_acknowledgement,
%       % text=black,
%       % font=\bfseries] {\insertdate};
%   \end{tikzpicture}
%   \setbeamercolor{title}{fg=title_page_title_colour}
%   \setbeamerfont{title}{size=\Large,series=\bfseries}
%   \setbeamercolor{author}{fg=title_page_author_colour}
%   \setbeamerfont{author}{size=\large,series=\bfseries}
%   \setbeamercolor{institute}{fg=title_page_institute_colour}
%   \setbeamerfont{institute}{size=\large,series=\bfseries}
%   \setbeamercolor{date}{fg=title_page_date_colour}
%   \setbeamerfont{date}{series=\bfseries}
% 	\titlepage
% \end{frame}
% \addtocounter{page}{-1}
% }

\newcommand{\titlepagewithfigure}[1]{%
  \begin{frame}[noframenumbering,plain]
    \begin{tikzpicture}[remember picture,overlay]
      \node[above right,inner sep=0pt,opacity=0.1] at (current page.south west)
      {
          \includegraphics[height=\paperheight,width=\paperwidth]{#1}
      };
      \node[anchor=north east,
      inner sep=5pt,
      opacity=0.9] at (current page.north east)
      {
          \includegraphics[width=0.2\textwidth]{FIGS-slides-admin/UM-logo-horizontal-CMYK.png}
      };
      \node[anchor=south, 
      align=justify, 
      text=black, 
      text width=1.1\textwidth,
      font=\footnotesize]  (land_acknowledgement)
      at (current page.south) 
      {The University of Manitoba campuses are located on original lands of Anishinaabeg, Ininew, Anisininew, Dakota and Dene peoples, and on the National Homeland of the Red River Métis.
      We respect the Treaties that were made on these territories, we acknowledge the harms and mistakes of the past, and we dedicate ourselves to move forward in partnership with Indigenous communities in a spirit of Reconciliation and collaboration.};  
      % \node[align=center, anchor=south,
      % above=0.5cm of land_acknowledgement,
      % text=black,
      % font=\bfseries] {\insertdate};
  \end{tikzpicture}
  \setbeamercolor{title}{fg=title_page_title_colour}
  \setbeamerfont{title}{size=\Large,series=\bfseries,family=\usefont{T1}{phv}{b}{n}}
  \setbeamercolor{author}{fg=title_page_author_colour}
  \setbeamerfont{author}{size=\large,series=\bfseries,family=\usefont{T1}{phv}{b}{n}}
  \setbeamercolor{institute}{fg=title_page_institute_colour}
  \setbeamerfont{institute}{size=\large,series=\bfseries,family=\usefont{T1}{phv}{b}{n}}
  \setbeamercolor{date}{fg=title_page_date_colour}
  \setbeamerfont{date}{series=\bfseries,family=\usefont{T1}{phv}{b}{n}}
	\titlepage
\end{frame}
\addtocounter{page}{-1}
}
% The outline page, with figure
% \newcommand{\outlinepage}[1]{%
% \begin{frame}[noframenumbering,plain]
%   \begin{tikzpicture}[remember picture,overlay]
%     \node[above right,inner sep=0pt,opacity=0.2] at (current page.south west)
%     {
%         \includegraphics[height=\paperheight,width=\paperwidth]{#1}
%     };
%   \end{tikzpicture}
%   \setbeamercolor{section in toc}{fg=outline_page_list_colour}
%   \setbeamerfont{section in toc}{size=\Large,series=\bfseries,family=\sffamily}
%   \frametitle{\textcolor{outline_page_title_colour}{\LARGE\bfseries Outline}}
%   \tableofcontents[hideallsubsections]
% \end{frame}
% \addtocounter{page}{-1}
% }
% The outline page, with figure
\newcommand{\outlinepage}[1]{%
\begin{frame}[noframenumbering,plain]
  \begin{tikzpicture}[remember picture,overlay]
    \node[above right,inner sep=0pt,opacity=0.2] at (current page.south west)
    {
        \includegraphics[height=\paperheight,width=\paperwidth]{#1}
    };
  \end{tikzpicture}
  \setbeamercolor{section in toc}{fg=outline_page_list_colour}
  % Use Helvetica Bold only for the outline slide TOC
  \setbeamerfont{section in toc}{size=\Large,family=\usefont{T1}{phv}{b}{n}}
  % Use Helvetica Bold for the outline title
  \frametitle{\textcolor{outline_page_title_colour}{\usefont{T1}{phv}{b}{n}\LARGE Outline}}
  \tableofcontents[hideallsubsections]
\end{frame}
\addtocounter{page}{-1}
}


%\let\oldsection\section
%\renewcommand{\section}[2]{\oldsection[#1]\newSectionSlide[#2]}


%%%%%%%%%%%%%%%%%%%%%
% CUSTOM SLIDE BACKGROUNDS
%%%%%%%%%%%%%%%%%%%%%
% Define custom background templates for different colors
\defbeamertemplate*{background canvas}{blue}{%
  \color{blue!15}\vrule width\paperwidth height\paperheight%
}
\defbeamertemplate*{background canvas}{green}{%
  \color{green!15}\vrule width\paperwidth height\paperheight%
}
\defbeamertemplate*{background canvas}{red}{%
  \color{red!15}\vrule width\paperwidth height\paperheight%
}
\defbeamertemplate*{background canvas}{yellow}{%
  \color{yellow!20}\vrule width\paperwidth height\paperheight%
}
\defbeamertemplate*{background canvas}{purple}{%
  \color{purple!15}\vrule width\paperwidth height\paperheight%
}
\defbeamertemplate*{background canvas}{orange}{%
  \color{orange!20}\vrule width\paperwidth height\paperheight%
}

% Define keys for the different background options
\makeatletter
\define@key{beamerframe}{blue}[true]{\setbeamertemplate{background canvas}[blue]}
\define@key{beamerframe}{green}[true]{\setbeamertemplate{background canvas}[green]}
\define@key{beamerframe}{red}[true]{\setbeamertemplate{background canvas}[red]}
\define@key{beamerframe}{yellow}[true]{\setbeamertemplate{background canvas}[yellow]}
\define@key{beamerframe}{purple}[true]{\setbeamertemplate{background canvas}[purple]}
\define@key{beamerframe}{orange}[true]{\setbeamertemplate{background canvas}[orange]}
\makeatother

% Reset to normal background for all frames by default
\BeforeBeginEnvironment{frame}{\setbeamertemplate{background canvas}[default]}

\newcommand{\punnett}[2]{
    \begin{center}
    \renewcommand{\arraystretch}{1.5} % Add space to rows
    \begin{tabular}{|c c | c|c|}
        \multicolumn{2}{c}{} & \multicolumn{2}{c}{\textbf{Father}} \\
        \multicolumn{2}{c}{} & #1 \\ \hline
        #2
    \end{tabular}
    \end{center}
}

% Colour definitions for Punnett squares
\usepackage{multirow}
\usepackage{colortbl}
\colorlet{punnettorange}{orange!30}
\colorlet{punnettblack}{black!50}
\colorlet{punnetttortie}{orange!60!black}

\usetikzlibrary{positioning, decorations.pathreplacing, arrows.meta}

% --- TIKZ STYLES FOR CAT DIAGRAMS ---
\tikzset{
    % This is the main placeholder style for the cat images
    catnode/.style={
        draw, 
        rectangle, 
        rounded corners, 
        minimum height=1.5cm, 
        minimum width=2.5cm, 
        align=center, 
        font=\small\bfseries
    },
    % Allele styles (the X and Y at the top/side)
    allele/.style={font=\Large\bfseries},
    female/.style={allele, color=purple!80!black},
    male/.style={allele, color=blue!80!black},
    % --- Placeholder styles for different cats ---
    % YOU CAN REPLACE THE CONTENTS OF THESE NODES WITH YOUR IMAGES
    orange_cat/.style={
        catnode, 
        fill=orange!30, 
        text=black
    },
    black_cat/.style={
        catnode, 
        fill=black!70, 
        text=white
    },
    tortie_cat/.style={
        catnode, 
        fill=orange!50!black, % A mix for tortoiseshell
        text=white
    }
}

\usetikzlibrary{
    arrows.meta, % For nicer arrow heads (e.g., -Latex)
    positioning, % For relative node placement (e.g., above=of)
    automata     % For state diagrams, loops
}
% --- END TIKZ STYLES ---

% Common date for all slides
\date{Fall 2026}
% Counter for random transition slides generated by R
\newcounter{sssectionbg}
\setcounter{sssectionbg}{1}
\newcommand{\SSsection}[1]{
  \Ssection{#1}{FIGS/random_bg_\thesssectionbg.pdf}
  \stepcounter{sssectionbg}
}

% Random transition slide generated by R for the title page
\newcommand{\SStitlepage}{
  \titlepagewithfigure{FIGS/random_title_bg.pdf}
}

% Random transition slide generated by R for the outline page
\newcommand{\SSoutlinepage}{
  \outlinepage{FIGS/random_outline_bg.pdf}
}
\newcommand{\SSsubsection}[1]{
  \Ssubsection{#1}{FIGS/random_bg_\thesssectionbg.pdf}
  \stepcounter{sssectionbg}
}

\newcommand{\SSsubsubsection}[1]{
  \Ssubsubsection{#1}{FIGS/random_bg_\thesssectionbg.pdf}
  \stepcounter{sssectionbg}
}



\usecolortheme{orchid}
%% Listings
\usepackage{listings}
\definecolor{mygreen}{rgb}{0,0.6,0}
\definecolor{mygray}{rgb}{0.5,0.5,0.5}
\definecolor{mymauve}{rgb}{0.58,0,0.82}
\definecolor{mygold}{rgb}{1,0.843,0}
\definecolor{myblue}{rgb}{0.537,0.812,0.941}

\definecolor{mygold2}{RGB}{120,105,22}
\definecolor{mygrey2}{RGB}{50,50,50}

\definecolor{lgreen}{rgb}{0.6,0.9,.6}
\definecolor{lred}{rgb}{1,0.5,.5}

\lstloadlanguages{R}
\lstset{ %
  language=R,
  backgroundcolor=\color{black!05},   % choose the background color
  basicstyle=\footnotesize\ttfamily,        % size of fonts used for the code
  breaklines=true,                 % automatic line breaking only at whitespace
  captionpos=b,                    % sets the caption-position to bottom
  commentstyle=\color{mygreen},    % comment style
  escapeinside={\%*}{*)},          % if you want to add LaTeX within your code
  keywordstyle=\color{red},       % keyword style
  stringstyle=\color{mygold},     % string literal style
  keepspaces=true,
  columns=fullflexible,
  tabsize=4,
}
% Could also do (in lstset)
% basicstyle==\fontfamily{pcr}\footnotesize
\lstdefinelanguage{Renhanced}%
  {keywords={abbreviate,abline,abs,acos,acosh,action,add1,add,%
      aggregate,alias,Alias,alist,all,anova,any,aov,aperm,append,apply,%
      approx,approxfun,apropos,Arg,args,array,arrows,as,asin,asinh,%
      atan,atan2,atanh,attach,attr,attributes,autoload,autoloader,ave,%
      axis,backsolve,barplot,basename,besselI,besselJ,besselK,besselY,%
      beta,binomial,body,box,boxplot,break,browser,bug,builtins,bxp,by,%
      c,C,call,Call,case,cat,category,cbind,ceiling,character,char,%
      charmatch,check,chol,chol2inv,choose,chull,class,close,cm,codes,%
      coef,coefficients,co,col,colnames,colors,colours,commandArgs,%
      comment,complete,complex,conflicts,Conj,contents,contour,%
      contrasts,contr,control,helmert,contrib,convolve,cooks,coords,%
      distance,coplot,cor,cos,cosh,count,fields,cov,covratio,wt,CRAN,%
      create,crossprod,cummax,cummin,cumprod,cumsum,curve,cut,cycle,D,%
      data,dataentry,date,dbeta,dbinom,dcauchy,dchisq,de,debug,%
      debugger,Defunct,default,delay,delete,deltat,demo,de,density,%
      deparse,dependencies,Deprecated,deriv,description,detach,%
      dev2bitmap,dev,cur,deviance,off,prev,,dexp,df,dfbetas,dffits,%
      dgamma,dgeom,dget,dhyper,diag,diff,digamma,dim,dimnames,dir,%
      dirname,dlnorm,dlogis,dnbinom,dnchisq,dnorm,do,dotplot,double,%
      download,dpois,dput,drop,drop1,dsignrank,dt,dummy,dump,dunif,%
      duplicated,dweibull,dwilcox,dyn,edit,eff,effects,eigen,else,%
      emacs,end,environment,env,erase,eval,equal,evalq,example,exists,%
      exit,exp,expand,expression,External,extract,extractAIC,factor,%
      fail,family,fft,file,filled,find,fitted,fivenum,fix,floor,for,%
      For,formals,format,formatC,formula,Fortran,forwardsolve,frame,%
      frequency,ftable,ftable2table,function,gamma,Gamma,gammaCody,%
      gaussian,gc,gcinfo,gctorture,get,getenv,geterrmessage,getOption,%
      getwd,gl,glm,globalenv,gnome,GNOME,graphics,gray,grep,grey,grid,%
      gsub,hasTsp,hat,heat,help,hist,home,hsv,httpclient,I,identify,if,%
      ifelse,Im,image,\%in\%,index,influence,measures,inherits,install,%
      installed,integer,interaction,interactive,Internal,intersect,%
      inverse,invisible,IQR,is,jitter,kappa,kronecker,labels,lapply,%
      layout,lbeta,lchoose,lcm,legend,length,levels,lgamma,library,%
      licence,license,lines,list,lm,load,local,locator,log,log10,log1p,%
      log2,logical,loglin,lower,lowess,ls,lsfit,lsf,ls,machine,Machine,%
      mad,mahalanobis,make,link,margin,match,Math,matlines,mat,matplot,%
      matpoints,matrix,max,mean,median,memory,menu,merge,methods,min,%
      missing,Mod,mode,model,response,mosaicplot,mtext,mvfft,na,nan,%
      names,omit,nargs,nchar,ncol,NCOL,new,next,NextMethod,nextn,%
      nlevels,nlm,noquote,NotYetImplemented,NotYetUsed,nrow,NROW,null,%
      numeric,\%o\%,objects,offset,old,on,Ops,optim,optimise,optimize,%
      options,or,order,ordered,outer,package,packages,page,pairlist,%
      pairs,palette,panel,par,parent,parse,paste,path,pbeta,pbinom,%
      pcauchy,pchisq,pentagamma,persp,pexp,pf,pgamma,pgeom,phyper,pico,%
      pictex,piechart,Platform,plnorm,plogis,plot,pmatch,pmax,pmin,%
      pnbinom,pnchisq,pnorm,points,poisson,poly,polygon,polyroot,pos,%
      postscript,power,ppoints,ppois,predict,preplot,pretty,Primitive,%
      print,prmatrix,proc,prod,profile,proj,prompt,prop,provide,%
      psignrank,ps,pt,ptukey,punif,pweibull,pwilcox,q,qbeta,qbinom,%
      qcauchy,qchisq,qexp,qf,qgamma,qgeom,qhyper,qlnorm,qlogis,qnbinom,%
      qnchisq,qnorm,qpois,qqline,qqnorm,qqplot,qr,Q,qty,qy,qsignrank,%
      qt,qtukey,quantile,quasi,quit,qunif,quote,qweibull,qwilcox,%
      rainbow,range,rank,rbeta,rbind,rbinom,rcauchy,rchisq,Re,read,csv,%
      csv2,fwf,readline,socket,real,Recall,rect,reformulate,regexpr,%
      relevel,remove,rep,repeat,replace,replications,report,require,%
      resid,residuals,restart,return,rev,rexp,rf,rgamma,rgb,rgeom,R,%
      rhyper,rle,rlnorm,rlogis,rm,rnbinom,RNGkind,rnorm,round,row,%
      rownames,rowsum,rpois,rsignrank,rstandard,rstudent,rt,rug,runif,%
      rweibull,rwilcox,sample,sapply,save,scale,scan,scan,screen,sd,se,%
      search,searchpaths,segments,seq,sequence,setdiff,setequal,set,%
      setwd,show,sign,signif,sin,single,sinh,sink,solve,sort,source,%
      spline,splinefun,split,sqrt,stars,start,stat,stem,step,stop,%
      storage,strstrheight,stripplot,strsplit,structure,strwidth,sub,%
      subset,substitute,substr,substring,sum,summary,sunflowerplot,svd,%
      sweep,switch,symbol,symbols,symnum,sys,status,system,t,table,%
      tabulate,tan,tanh,tapply,tempfile,terms,terrain,tetragamma,text,%
      time,title,topo,trace,traceback,transform,tri,trigamma,trunc,try,%
      ts,tsp,typeof,unclass,undebug,undoc,union,unique,uniroot,unix,%
      unlink,unlist,unname,untrace,update,upper,url,UseMethod,var,%
      variable,vector,Version,vi,warning,warnings,weighted,weights,%
      which,while,window,write,\%x\%,x11,X11,xedit,xemacs,xinch,xor,%
      xpdrows,xy,xyinch,yinch,zapsmall,zip},%
   otherkeywords={!,!=,~,$,*,\%,\&,\%/\%,\%*\%,\%\%,<-,<<-,_,/},%
   alsoother={._$},%
   sensitive,%
   morecomment=[l]\#,%
   morestring=[d]",%
   morestring=[d]'% 2001 Robert Denham
  }%

% Custom color for listings
\definecolor{codegreen}{rgb}{0,0.6,0}
\definecolor{codegray}{rgb}{0.5,0.5,0.5}
\definecolor{codepurple}{rgb}{0.58,0,0.82}
\definecolor{backcolour}{rgb}{0.95,0.95,0.92}

%------------------------------------------------------------
% Listings configuration for Julia
%------------------------------------------------------------
% Define the Julia language for the listings package
\lstdefinelanguage{Julia}
{
  morekeywords={abstract,baremodule,begin,break,catch,const,continue,do,else,elseif,end,export,false,finally,for,function,global,if,import,let,local,macro,module,quote,return,struct,true,try,using,while,in},
  sensitive=true,
  comment=[l]{\#},
  morecomment=[n]{\#=}{=\#}, % For block comments like #= ... =#
  string=[b]",
  morestring=[b]',
  alsoletter={!}, % Allow ! in identifiers for mutating functions
}

% Define the style for the code blocks
\lstdefinestyle{JuliaStyle}{
    backgroundcolor=\color{backcolour},   
    commentstyle=\color{codegreen},
    keywordstyle=\color{magenta},
    numberstyle=\tiny\color{codegray},
    stringstyle=\color{codepurple},
    basicstyle=\ttfamily\footnotesize,
    breakatwhitespace=false,         
    breaklines=true,                 
    captionpos=b,                    
    keepspaces=true,                 
    % numbers=left,                    
    % numbersep=5pt,                  
    showspaces=false,                
    showstringspaces=false,
    showtabs=false,                  
    tabsize=2,
    language=Julia % Use the newly defined Julia language
}
\lstset{style=JuliaStyle}

%%%%%%% 
%% Definitions in yellow boxes
\usepackage{etoolbox}
\setbeamercolor{block title}{use=structure,fg=structure.fg,bg=structure.fg!40!bg}
\setbeamercolor{block body}{parent=normal text,use=block title,bg=block title.bg!20!bg}

% Important definition
\BeforeBeginEnvironment{importantdefinition}{%
	\setbeamercolor{block title}{fg=black,bg=yellow!20!white}
	\setbeamercolor{block body}{fg=black, bg=yellow!05!white}
}
\AfterEndEnvironment{importantdefinition}{
	\setbeamercolor{block title}{use=structure,fg=structure.fg,bg=structure.fg!20!bg}
	\setbeamercolor{block body}{parent=normal text,use=block title,bg=block title.bg!50!bg, fg=black}
}
% Regular (less important) definition
\BeforeBeginEnvironment{definition}{%
	\setbeamercolor{block title}{fg=black,bg=black!15!white}
	\setbeamercolor{block body}{fg=black, bg=black!05!white}
}
\AfterEndEnvironment{definition}{
	\setbeamercolor{block title}{use=structure,fg=structure.fg,bg=structure.fg!20!bg}
	\setbeamercolor{block body}{parent=normal text,use=block title,bg=block title.bg!50!bg, fg=black}
}
% Important theorem (red)
\BeforeBeginEnvironment{importanttheorem}{%
	\setbeamercolor{block title}{fg=black,bg=red!20!white}
	\setbeamercolor{block body}{fg=black, bg=red!05!white}
}
\AfterEndEnvironment{importanttheorem}{
	\setbeamercolor{block title}{use=structure,fg=structure.fg,bg=structure.fg!20!bg}
	\setbeamercolor{block body}{parent=normal text,use=block title,bg=block title.bg!50!bg, fg=black}
}
\BeforeBeginEnvironment{importantproperty}{%
	\setbeamercolor{block title}{fg=black,bg=red!50!white}
	\setbeamercolor{block body}{fg=black, bg=red!30!white}
}
\AfterEndEnvironment{importantproperty}{
	\setbeamercolor{block title}{use=structure,fg=structure.fg,bg=structure.fg!20!bg}
	\setbeamercolor{block body}{parent=normal text,use=block title,bg=block title.bg!50!bg, fg=black}
}

%% Colours for title page, outline and the like
\definecolor{UMbrown}{RGB}{79,44,29}
\definecolor{HighlightGold}{RGB}{230,165,83}
\definecolor{DeepGreen}{RGB}{27,46,28}
\definecolor{MediumGreen}{RGB}{52,128,56}
\definecolor{LightGreen}{RGB}{150,224,154}
\definecolor{AccentRed}{RGB}{196,44,27}
%% These are the titles used in practice
% Title page
\colorlet{title_page_title_colour}{UMbrown}
\colorlet{title_page_author_colour}{UMbrown}
\colorlet{title_page_date_colour}{UMbrown}
\colorlet{title_page_institute_colour}{UMbrown}
% Outline page
\colorlet{outline_page_title_colour}{DeepGreen}
\colorlet{outline_page_list_colour}{UMbrown}
% Section page
\colorlet{section_page_title_colour}{DeepGreen}
\colorlet{section_page_list_colour}{UMbrown}
% Subection page
\colorlet{subsection_page_title_colour}{DeepGreen}
\colorlet{subsection_page_list_colour}{UMbrown}





\IfFileExists{upquote.sty}{\usepackage{upquote}}{}
\begin{document}

% Set up cross-references and counter persistence
\setcounter{theorem}{166}

%%%%%%%%%%%%%%%%%%%%%%%%%%%%%%%%%
%%%%%%%%%%%%%%%%%%%%%%%%%%%%%%%%%
%% TITLE AND OUTLINE
%%%%%%%%%%%%%%%%%%%%%%%%%%%%%%%%%
%%%%%%%%%%%%%%%%%%%%%%%%%%%%%%%%%
\SStitlepage
\SSoutlinepage


\begin{frame}{Why represent text using linear algebra?}
\begin{itemize}
\item Much of the data generated in science, industry and daily life consists of unstructured text: books, articles, clinical notes and web pages
\vfill
\item Computers cannot perform calculus or geometry directly on strings of characters
\vfill
\item In this unit, we turn discrete text corpora into \defword{sparse non-negative matrices} and apply three core tools of linear algebra:
\begin{enumerate}
\item \textbf{Diagonal matrix scaling} (TF-IDF) to reweight rows and columns
\item \textbf{Inner products and angles} (cosine similarity) to measure semantic similarity in high dimensions
\item \textbf{Constrained low-rank factorisation} (Non-negative Matrix Factorisation via Alternating Least Squares) to discover latent topics
\end{enumerate}
\end{itemize}
\end{frame}


%%%%%%%%%%%%%%%%%%%%%
%%%%%%%%%%%%%%%%%%%%%
%%%%%%%%%%%%%%%%%%%%%
%%%%%%%%%%%%%%%%%%%%%
\SSsection{Text tokenisation and the vector space model}

\begin{frame}{The central idea: from words to vectors}
\begin{itemize}
\item The fundamental challenge in natural language processing (\defword{NLP}) and information retrieval is to translate discrete symbolic data (\defword{words} and \defword{documents}) into a numerical format
\vfill
\item Why? Mathematical models are built on vector spaces, inner products and matrix factorisations rather than dictionary definitions
\vfill
\item The \defword{vector space model} (\defword{VSM}), introduced in information retrieval by Salton, Wong and Yang (1975), represents documents and words as vectors in $\IR^m$ so that the geometry of the space reflects semantic relationships
\vfill
\item Documents or words with similar meaning are mapped to vectors pointing in similar directions in $\IR^m$
\end{itemize}
\end{frame}


\begin{frame}{A visual metaphor for the vector space model}
\begin{center}
\begin{tikzpicture}[
    scale=0.85,
    dot/.style={circle, fill=blue!60!black, inner sep=2pt},
    lbl/.style={font=\small}
]
    % Axes
    \draw[->, thick] (0,0) -- (6,0) node[right, font=\small\bfseries] {Dimension 1 (algebra)};
    \draw[->, thick] (0,0) -- (0,4.5) node[above, font=\small\bfseries] {Dimension 2 (graphs)};

    % Vectors
    \draw[->, thick, blue!70!black] (0,0) -- (4.8,0.9);
    \draw[->, thick, blue!70!black] (0,0) -- (4.2,1.2);
    \draw[->, thick, orange!85!black] (0,0) -- (0.9,3.9);
    \draw[->, thick, orange!85!black] (0,0) -- (1.3,3.5);

    % Points and labels
    \node[dot] at (4.8,0.9) {};
    \node[lbl, right] at (4.8,0.9) {doc 1 (matrices)};
    \node[dot] at (4.2,1.2) {};
    \node[lbl, above left] at (4.2,1.2) {doc 2 (eigenvalues)};
    \node[dot] at (0.9,3.9) {};
    \node[lbl, right] at (0.9,3.9) {doc 3 (networks)};
    \node[dot] at (1.3,3.5) {};
    \node[lbl, right] at (1.3,3.5) {doc 4 (shortest paths)};
\end{tikzpicture}
\end{center}
\vfill
\begin{itemize}
\item Vectors for documents on linear algebra point in nearly the same direction
\item Vectors for documents on graph theory cluster along a different direction
\item The angle between a linear algebra document and a graph theory document is large
\end{itemize}
\end{frame}


\begin{frame}{What is a token?}
\begin{itemize}
\item A \defword{token} is the fundamental discrete unit of text processed by a model: it can be a word, a word stem, a sub-word or a character
\vfill
\item Converting a raw character string into a sequence of tokens is called \defword{tokenisation}
\vfill
\item For example, after lower-casing and punctuation removal, the sentence ``The cat sat on the mat'' is tokenised as
\[
[\code{"the"},\; \code{"cat"},\; \code{"sat"},\; \code{"on"},\; \code{"the"},\; \code{"mat"}]
\]
\vfill
\item Each unique token in our collection is assigned an integer index $i \in \{1,\ldots,m\}$ from a fixed vocabulary, yielding an integer sequence such as $[1, 2, 3, 4, 1, 5]$
\end{itemize}
\end{frame}


\begin{frame}{Vocabulary and standard basis (one-hot) vectors}
\begin{importantdefinition}[Vocabulary and tokenisation]
A \defword{vocabulary} $V = \{w_1, w_2, \ldots, w_m\}$ is an ordered set of $m = |V|$ distinct tokens. \defword{Tokenisation} maps a raw document $d$ to a finite sequence of tokens $(t_1, t_2, \ldots, t_L)$ with each $t_\ell \in V$
\end{importantdefinition}
\vfill
\begin{importantdefinition}[One-hot representation of a token]
Given a vocabulary $V = \{w_1, \ldots, w_m\}$ of size $m$, the \defword{one-hot vector} associated with token $w_i \in V$ is the $i$-th standard basis vector $\be_i \in \IR^m$, whose $k$-th entry is
\[
(\be_i)_k = \delta_{ik} =
\begin{cases}
1 & \textrm{if } k = i \\
0 & \textrm{if } k \neq i
\end{cases}
\]
\end{importantdefinition}
\end{frame}


\begin{frame}{Visualising one-hot vectors and their limitations}
\begin{columns}[c]
\begin{column}{0.45\textwidth}
\begin{center}
\begin{tikzpicture}[scale=0.9]
    \matrix [matrix of nodes, nodes in empty cells, draw, thick, column sep=0mm, row sep=0mm,
             cells={minimum height=0.6cm, minimum width=1.2cm, align=center}] (vector) {
             $0$ \\ $0$ \\ $1$ \\ $0$ \\ $\vdots$ \\ $0$ \\
    };
    \node[left, xshift=-0.3cm, font=\small] at (vector-1-1.west) {Index 1: \code{"cat"}};
    \node[left, xshift=-0.3cm, font=\small] at (vector-2-1.west) {Index 2: \code{"dog"}};
    \node[left, xshift=-0.3cm, font=\small\bfseries, color=orange!90!black] at (vector-3-1.west) {Index 3: \code{"matrix"}};
    \node[left, xshift=-0.3cm, font=\small] at (vector-4-1.west) {Index 4: \code{"vector"}};
    \node[left, xshift=-0.3cm, font=\small] at (vector-6-1.west) {Index $m$: \code{"zebra"}};
    \node[right, xshift=0.3cm, font=\normalsize] at (vector.east) {$= \be_3 \in \IR^m$};
\end{tikzpicture}
\end{center}
\end{column}
\begin{column}{0.55\textwidth}
\begin{itemize}
\item The set $\{\be_1, \ldots, \be_m\}$ is the standard orthonormal basis of $\IR^m$
\vfill
\item For any two distinct tokens $w_i \neq w_j$:
\[
\be_i \bullet \be_j = 0 \quad\textrm{and}\quad \|\be_i - \be_j\|_2 = \sqrt{2}
\]
\vfill
\item One-hot vectors alone treat every pair of distinct words as equally unrelated: \code{"matrix"} is as far from \code{"vector"} as it is from \code{"cat"}
\end{itemize}
\end{column}
\end{columns}
\end{frame}


\begin{frame}{From tokens to documents: the bag-of-words model}
How do we represent a whole document $d_j$ as a single vector in $\IR^m$?
\vfill
\begin{importantdefinition}[Bag-of-words document vector]
Let $V = \{w_1, \ldots, w_m\}$ be a vocabulary of size $m$ and let $d_j$ be a tokenised document $(t_1, t_2, \ldots, t_{L_j})$. The \defword{bag-of-words} (\defword{BoW}) vector of document $d_j$ is
\[
\bx_j = \sum_{\ell=1}^{L_j} \be_{\mathsf{idx}(t_\ell)} = \sum_{i=1}^m f_{ij} \be_i =
\begin{pmatrix}
f_{1j} \\ f_{2j} \\ \vdots \\ f_{mj}
\end{pmatrix} \in \IRplus^m
\]
where $f_{ij} \in \IN_0$ is the number of times term $w_i$ occurs in document $d_j$
\end{importantdefinition}
\end{frame}


\begin{frame}{Example -- Constructing bag-of-words vectors}
Suppose our vocabulary has $m = 5$ terms:
\[
V = \{w_1=\code{"cat"},\; w_2=\code{"dog"},\; w_3=\code{"mat"},\; w_4=\code{"sat"},\; w_5=\code{"the"}\}
\]
\vfill
Consider two short documents (after lower-casing):
\begin{itemize}
\item $d_1$: \code{"the cat sat on the mat"} (ignoring \code{"on"} $\notin V$)
\item $d_2$: \code{"the dog sat on the dog mat"}
\end{itemize}
\vfill
Summing the standard basis vectors gives their bag-of-words vectors in $\IRplus^5$:
\[
\bx_1 = 1\be_1 + 0\be_2 + 1\be_3 + 1\be_4 + 2\be_5 = \begin{pmatrix} 1 \\ 0 \\ 1 \\ 1 \\ 2 \end{pmatrix}, \qquad
\bx_2 = 0\be_1 + 2\be_2 + 1\be_3 + 1\be_4 + 2\be_5 = \begin{pmatrix} 0 \\ 2 \\ 1 \\ 1 \\ 2 \end{pmatrix}
\]
\end{frame}


\begin{frame}{What bag-of-words preserves and what it discards}
\begin{itemize}
\item Because vector addition in $\IR^m$ is commutative ($\be_i + \be_k = \be_k + \be_i$), the bag-of-words map completely discards \textbf{word order} and grammar
\vfill
\item For example, ``dog bites man'' and ``man bites dog'' produce the exact same bag-of-words vector $\bx \in \IRplus^m$
\vfill
\item Despite this simplification, bag-of-words vectors are remarkably effective for \textbf{topic classification}, \textbf{document clustering} and \textbf{search}:
\begin{itemize}
\item Documents about linear algebra frequently contain \code{"matrix"}, \code{"rank"} and \code{"eigenvalue"}
\item Documents about graph theory frequently contain \code{"vertex"}, \code{"edge"} and \code{"path"}
\end{itemize}
\end{itemize}
\end{frame}


\begin{frame}{Sparsity and non-negativity in high dimensions}
\begin{definition}[Support and number of non-zero entries]
For a vector $\bx = (x_1, \ldots, x_m)^T \in \IR^m$, its \defword{support} is the index set of non-zero entries
\[
\mathsf{supp}(\bx) = \{i \in \{1, \ldots, m\} : x_i \neq 0\}
\]
and we write $\|\bx\|_0 = |\mathsf{supp}(\bx)|$ for the number of non-zero entries of $\bx$
\end{definition}
\vfill
\begin{itemize}
\item Note that $\|\bx\|_0$ is \emph{not} a true norm (since $\|\alpha \bx\|_0 = \|\bx\|_0$ for $\alpha \neq 0$), but it measures \defword{sparsity}
\vfill
\item In realistic text corpora, the vocabulary size $m$ is large ($10^4$ to $10^5$ terms), whereas a single document uses only a few hundred distinct words
\vfill
\item Thus $\|\bx_j\|_0 \ll m$: every document vector $\bx_j \in \IRplus^m$ is a \textbf{high-dimensional, sparse, non-negative vector} ($98\%$ to $99.9\%$ of its entries are zero)
\end{itemize}
\end{frame}


\begin{frame}{Vocabulary pre-processing as row selection}
In practice, before forming document vectors $\bx_j \in \IRplus^m$, raw tokens undergo two standard pre-processing steps:
\vfill
\begin{enumerate}
\item \textbf{Stop-word removal}: discarding high-frequency function words (\code{"the"}, \code{"and"}, \code{"of"}, \code{"to"}) that carry little topical meaning
\item \textbf{Stemming or lemmatisation}: collapsing inflected forms (\code{"matrices"}, \code{"matrix"} $\to$ \code{"matrix"}; \code{"computing"}, \code{"computed"} $\to$ \code{"comput"}) into a single vocabulary row
\end{enumerate}
\vfill
\begin{itemize}
\item \textbf{Linear algebra view}: Removing $m - m'$ stop words from a vocabulary of size $m$ is equivalent to left-multiplying each document vector $\bx_j \in \IRplus^m$ by a \textbf{row-selection matrix} $P \in \M_{m', m}(\{0, 1\})$ obtained by deleting the stop-word rows of $\II_m$:
\[
\bx_j^{\text{reduced}} = P \bx_j \in \IRplus^{m'}
\]
\end{itemize}
\end{frame}


%%%%%%%%%%%%%%%%%%%%%
%%%%%%%%%%%%%%%%%%%%%
%%%%%%%%%%%%%%%%%%%%%
%%%%%%%%%%%%%%%%%%%%%
\SSsection{Term-document matrices and TF-IDF scaling}

%%%%%%%%%%%%%%%%%%%%%
%%%%%%%%%%%%%%%%%%%%%
\SSsubsection{The term-document matrix}

\begin{frame}{The term-document matrix}
\begin{importantdefinition}[Term-document matrix]
Given a vocabulary $V = \{w_1, \ldots, w_m\}$ of $m$ terms and a corpus $D = \{d_1, \ldots, d_n\}$ of $n$ documents, the \defword{term-document matrix} is the non-negative matrix
\[
X = [x_{ij}] = \begin{bmatrix} \bx_1 & \bx_2 & \cdots & \bx_n \end{bmatrix} \in \M_{m,n}(\IRplus)
\]
where:
\begin{itemize}
\item Row $i \in \{1, \ldots, m\}$ corresponds to vocabulary term $w_i$
\item Column $j \in \{1, \ldots, n\}$ is the vector $\bx_j \in \IRplus^m$ of document $d_j$
\item Entry $x_{ij} \ge 0$ is the weight (e.g., raw frequency $f_{ij}$) of term $w_i$ in document $d_j$
\end{itemize}
\end{importantdefinition}
\vfill
\begin{itemize}
\item Some software libraries work with the transpose $X^T \in \M_{n,m}(\IRplus)$, called the \defword{document-term matrix}; in linear algebra we prefer $X \in \M_{m,n}(\IRplus)$ so that documents are column vectors in $\IRplus^m$
\end{itemize}
\end{frame}


\begin{frame}{Example -- A toy corpus of $n=4$ documents and $m=6$ terms}
Consider $n=4$ short documents across two mathematical themes, tracked over the vocabulary $V = \{\code{matrix}, \code{vector}, \code{eigenvalue}, \code{graph}, \code{network}, \code{data}\}$ ($m=6$):
\vfill
\begin{itemize}
\item \textbf{Document $d_1$ (linear algebra)}: mentions \code{matrix} (3 times), \code{vector} (2 times), \code{eigenvalue} (2 times) and \code{data} (4 times)
\vfill
\item \textbf{Document $d_2$ (linear algebra)}: mentions \code{matrix} (2 times), \code{vector} (3 times) and \code{data} (3 times)
\vfill
\item \textbf{Document $d_3$ (graph theory)}: mentions \code{graph} (3 times), \code{network} (2 times) and \code{data} (5 times)
\vfill
\item \textbf{Document $d_4$ (graph theory)}: mentions \code{vector} (1 time), \code{graph} (2 times), \code{network} (3 times) and \code{data} (4 times)
\end{itemize}
\end{frame}


\begin{frame}{Example -- Explicit $6 \times 4$ term-document matrix $X$}
Assembling the document vectors $\bx_1, \bx_2, \bx_3, \bx_4 \in \IRplus^6$ as columns gives $X \in \M_{6,4}(\IRplus)$:
\[
\begin{array}{r|cccc}
\textrm{Term} & d_1 & d_2 & d_3 & d_4 \\
\hline
\code{matrix} & 3 & 2 & 0 & 0 \\
\code{vector} & 2 & 3 & 0 & 1 \\
\code{eigenvalue} & 2 & 0 & 0 & 0 \\
\code{graph} & 0 & 0 & 3 & 2 \\
\code{network} & 0 & 0 & 2 & 3 \\
\code{data} & 4 & 3 & 5 & 4
\end{array}
\qquad\implies\qquad
X = \begin{pmatrix}
3 & 2 & 0 & 0 \\
2 & 3 & 0 & 1 \\
2 & 0 & 0 & 0 \\
0 & 0 & 3 & 2 \\
0 & 0 & 2 & 3 \\
4 & 3 & 5 & 4
\end{pmatrix}
\]
\vfill
\begin{itemize}
\item Notice the near-block structure in rows $1$--$5$ separating $\{d_1, d_2\}$ from $\{d_3, d_4\}$
\item In contrast, row $6$ (\code{data}) has large entries across \textbf{all four columns}, obscuring the two topics!
\end{itemize}
\end{frame}


\begin{frame}{Dual geometric views of $X \in \M_{m,n}(\IRplus)$}
Every term-document matrix $X \in \M_{m,n}(\IRplus)$ admits two complementary geometric interpretations:
\vfill
\begin{columns}[t]
\begin{column}{0.48\textwidth}
\textbf{1. Column view (document space $\IRplus^m$)}
\begin{itemize}
\item Each column $\bx_j \in \IRplus^m$ ($j=1,\ldots,n$) represents a \textbf{document} as a point in $m$-dimensional term space
\item Used to compare, cluster or classify documents
\end{itemize}
\end{column}
\begin{column}{0.48\textwidth}
\textbf{2. Row view (term space $\IRplus^n$)}
\begin{itemize}
\item Each transposed row $\bt_i \in \IRplus^n$ ($i=1,\ldots,m$) represents a \textbf{term} via its occurrence profile across the $n$ documents
\item Two words that co-occur in the same documents have similar row vectors
\end{itemize}
\end{column}
\end{columns}
\end{frame}


%%%%%%%%%%%%%%%%%%%%%
%%%%%%%%%%%%%%%%%%%%%
\SSsubsection{TF-IDF as diagonal matrix multiplication}

\begin{frame}{Why raw term counts fail}
Using raw frequency counts $x_{ij} = f_{ij}$ directly in $X$ suffers from two major defects:
\vfill
\begin{enumerate}
\item \textbf{Common words dominate rare discriminative words}
\begin{itemize}
\item Words like \code{"the"}, \code{"is"} or \code{"data"} appear in almost every document with large counts, dominating Euclidean norms and inner products
\item Rare domain terms like \code{"eigenvalue"} or \code{"isomorphism"} occur fewer times but carry far more information about a document's topic
\end{itemize}
\vfill
\item \textbf{Document length distorts vector magnitude}
\begin{itemize}
\item A 20-page article and a 1-paragraph abstract on the exact same topic have proportional term frequencies, but the longer document has a much larger norm $\|\bx_j\|_2$
\end{itemize}
\end{enumerate}
\end{frame}


\begin{frame}{Term frequency and inverse document frequency (TF-IDF)}
\begin{importantdefinition}[Document frequency, IDF and TF-IDF]
Let $X = [f_{ij}] \in \M_{m,n}(\IRplus)$ be a raw term-document count matrix for $m$ terms and $n$ documents:
\begin{enumerate}
\item The \defword{term frequency} $\mathsf{tf}_{ij}$ is the raw count $f_{ij}$ (or a scaled variant such as $f_{ij} / \sum_{k=1}^m f_{kj}$)
\item The \defword{document frequency} $\mathsf{df}_i$ of term $w_i$ is the number of documents in which $w_i$ appears:
\[
\mathsf{df}_i = \left|\{j \in \{1, \ldots, n\} : f_{ij} > 0\}\right| \in \{1, \ldots, n\}
\]
\item The \defword{inverse document frequency} of term $w_i$ is
\[
\mathsf{idf}_i = \log\!\left(\frac{n}{\mathsf{df}_i}\right) \ge 0
\]
\item The \defword{TF-IDF weight} of term $w_i$ in document $d_j$ is $\mathsf{tfidf}_{ij} = \mathsf{tf}_{ij} \cdot \mathsf{idf}_i$
\end{enumerate}
\end{importantdefinition}
\end{frame}


\begin{frame}{How $\mathsf{idf}_i = \log(n / \mathsf{df}_i)$ reweights vocabulary terms}
Assume every vocabulary term $w_i \in V$ appears in at least one document ($1 \le \mathsf{df}_i \le n$):
\vfill
\begin{center}
\begin{tabular}{lccl}
\toprule
\textbf{Term behaviour} & $\mathsf{df}_i$ & $\mathsf{idf}_i = \log(n/\mathsf{df}_i)$ & \textbf{Effect on row $i$ of $X$} \\
\midrule
Appears in all documents & $n$ & $\log(1) = 0$ & Zeroed out completely \\
Appears in half of documents & $n/2$ & $\log(2) \approx 0.693$ & Moderately weighted \\
Appears in 1 document only & $1$ & $\log(n)$ (maximal) & Strongly amplified \\
\bottomrule
\end{tabular}
\end{center}
\vfill
\begin{itemize}
\item Note: throughout this course, $\log(\cdot)$ denotes the natural logarithm $\ln(\cdot)$ (changing the base of the logarithm only multiplies all $\mathsf{idf}_i$ by a positive constant)
\item In practice, when out-of-corpus queries might have $\mathsf{df}_i = 0$, smoothed variants such as $\log\!\left(\frac{n+1}{\mathsf{df}_i+1}\right) + 1$ are also used
\end{itemize}
\end{frame}


\begin{frame}{TF-IDF and length normalisation as diagonal matrix scaling}
\begin{importanttheorem}[Row and column scaling via diagonal matrices]\label{th:tfidf_diag_scaling}
Let $X = [x_{ij}] \in \M_{m,n}(\IRplus)$ be a term-document matrix with non-zero columns $\bx_1, \ldots, \bx_n$. Define the two diagonal matrices
\[
D_{\text{IDF}} = \diag(\mathsf{idf}_1, \ldots, \mathsf{idf}_m) \in \M_m(\IRplus), \qquad
D_{\text{len}}^{-1} = \diag\!\left(\frac{1}{\|\bx_1\|}, \ldots, \frac{1}{\|\bx_n\|}\right) \in \M_n(\IRplus)
\]
Then:
\begin{enumerate}
\item \textbf{Left-multiplication} $X_{\text{TF-IDF}} = D_{\text{IDF}} X$ scales the \textbf{rows} (terms) of $X$: entry $(i,j)$ is $\mathsf{idf}_i \, x_{ij}$
\item \textbf{Right-multiplication} $X D_{\text{len}}^{-1}$ scales the \textbf{columns} (documents) of $X$: entry $(i,j)$ is $x_{ij} / \|\bx_j\|$
\item The combined matrix $\tilde{X} = D_{\text{IDF}} X D_{\text{len}}^{-1} \in \M_{m,n}(\IRplus)$ has entries $\tilde{x}_{ij} = \frac{\mathsf{idf}_i \, x_{ij}}{\|\bx_j\|}$
\end{enumerate}
\end{importanttheorem}
\end{frame}


\begin{frame}[red]{Proof of Theorem~\ref{th:tfidf_diag_scaling}}
Let $D_{\text{IDF}} = [d_{ik}] \in \M_m$ with $d_{ii} = \mathsf{idf}_i$ and $d_{ik} = 0$ for $k \neq i$. By definition of matrix multiplication, the $(i,j)$ entry of $D_{\text{IDF}} X$ is
\[
(D_{\text{IDF}} X)_{ij} = \sum_{k=1}^m d_{ik} x_{kj} = d_{ii} x_{ij} = \mathsf{idf}_i \, x_{ij}
\]
Thus row $i$ of $X$ is multiplied by the scalar $\mathsf{idf}_i$
\vfill
Similarly, let $D_{\text{len}}^{-1} = [c_{\ell j}] \in \M_n$ with $c_{jj} = 1/\|\bx_j\|$ and $c_{\ell j} = 0$ for $\ell \neq j$. For any matrix $A = [a_{i\ell}] \in \M_{m,n}$, the $(i,j)$ entry of $A D_{\text{len}}^{-1}$ is
\[
(A D_{\text{len}}^{-1})_{ij} = \sum_{\ell=1}^n a_{i\ell} c_{\ell j} = a_{ij} c_{jj} = \frac{a_{ij}}{\|\bx_j\|}
\]
Setting $A = D_{\text{IDF}} X$ gives $\tilde{x}_{ij} = (D_{\text{IDF}} X D_{\text{len}}^{-1})_{ij} = \frac{\mathsf{idf}_i \, x_{ij}}{\|\bx_j\|}$ \hfill\qed
\end{frame}


\begin{frame}{Visualising $\tilde{X} = D_{\text{IDF}} X D_{\text{len}}^{-1}$}
\begin{center}
\begin{tikzpicture}[scale=0.88, every node/.style={font=\small}]
    % TilX
    \draw[thick, fill=gray!15] (0,0) rectangle (2.2,3.2);
    \node at (1.1,1.6) {$\tilde{X}$};
    \node[below] at (1.1,0) {$m \times n$};

    \node[font=\Large] at (2.7,1.6) {$=$};

    % D_IDF
    \draw[thick, fill=orange!20] (3.2,0) rectangle (6.4,3.2);
    \draw[ultra thick, orange!85!black] (3.2,3.2) -- (6.4,0);
    \node[above right] at (3.25,0.1) {$0$};
    \node[below left] at (6.35,3.1) {$0$};
    \node[fill=orange!20, inner sep=1pt] at (4.8,1.6) {$D_{\text{IDF}}$};
    \node[below] at (4.8,0) {$m \times m$ (scales rows)};

    % X
    \draw[thick, fill=blue!15] (6.8,0) rectangle (9.0,3.2);
    \node at (7.9,1.6) {$X$};
    \node[below] at (7.9,0) {$m \times n$};

    % D_len^-1
    \draw[thick, fill=green!20] (9.4,1.0) rectangle (11.6,3.2);
    \draw[ultra thick, green!60!black] (9.4,3.2) -- (11.6,1.0);
    \node[fill=green!20, inner sep=1pt] at (10.5,2.1) {$D_{\text{len}}^{-1}$};
    \node[below] at (10.5,1.0) {$n \times n$ (scales cols)};
\end{tikzpicture}
\end{center}
\vfill
\begin{itemize}
\item Notice the fundamental pattern from elementary linear algebra:
\begin{itemize}
\item \textbf{Left-multiplying} by a diagonal matrix scales the \textbf{rows} (the $m$ vocabulary terms)
\item \textbf{Right-multiplying} by a diagonal matrix scales the \textbf{columns} (the $n$ documents)
\end{itemize}
\end{itemize}
\end{frame}


\begin{frame}{Example -- Computing $D_{\text{IDF}} X$ by hand (Part 1)}
Consider a corpus of $n = 4$ documents over $m = 4$ terms with raw term-document matrix $X \in \M_{4,4}(\IRplus)$:
\[
\begin{array}{r|cccc}
\textrm{Term} & d_1 & d_2 & d_3 & d_4 \\
\hline
w_1\ (\code{matrix}) & 3 & 2 & 0 & 0 \\
w_2\ (\code{eigenvalue}) & 2 & 0 & 0 & 0 \\
w_3\ (\code{graph}) & 0 & 0 & 1 & 3 \\
w_4\ (\code{data}) & 5 & 4 & 6 & 5
\end{array}
\qquad\implies\qquad
X = \begin{pmatrix}
3 & 2 & 0 & 0 \\
2 & 0 & 0 & 0 \\
0 & 0 & 1 & 3 \\
5 & 4 & 6 & 5
\end{pmatrix}
\]
\vfill
\textbf{Step 1: Count non-zero entries in each row to find $\mathsf{df}_i$ and $\mathsf{idf}_i = \log(4/\mathsf{df}_i)$}:
\begin{itemize}
\item Row 1 ($w_1$): appears in $d_1, d_2 \implies \mathsf{df}_1 = 2 \implies \mathsf{idf}_1 = \log(4/2) = \log 2$
\item Row 2 ($w_2$): appears in $d_1$ only $\implies \mathsf{df}_2 = 1 \implies \mathsf{idf}_2 = \log(4/1) = \log 4 = 2\log 2$
\item Row 3 ($w_3$): appears in $d_3, d_4 \implies \mathsf{df}_3 = 2 \implies \mathsf{idf}_3 = \log(4/2) = \log 2$
\item Row 4 ($w_4$): appears in $d_1, d_2, d_3, d_4 \implies \mathsf{df}_4 = 4 \implies \mathsf{idf}_4 = \log(4/4) = \log 1 = 0$
\end{itemize}
\end{frame}


\begin{frame}{Example -- Computing $D_{\text{IDF}} X$ by hand (Part 2)}
\textbf{Step 2: Form $D_{\text{IDF}}$ and left-multiply $X$}:
\[
D_{\text{IDF}} = \diag(\log 2,\; 2\log 2,\; \log 2,\; 0) = (\log 2)\begin{pmatrix} 1 & 0 & 0 & 0 \\ 0 & 2 & 0 & 0 \\ 0 & 0 & 1 & 0 \\ 0 & 0 & 0 & 0 \end{pmatrix}
\]
\vfill
Multiplying each row of $X$ by its diagonal entry gives
\[
X_{\text{TF-IDF}} = D_{\text{IDF}} X = (\log 2)
\begin{pmatrix}
3 & 2 & 0 & 0 \\
4 & 0 & 0 & 0 \\
0 & 0 & 1 & 3 \\
0 & 0 & 0 & 0
\end{pmatrix}
\approx
\begin{pmatrix}
2.08 & 1.39 & 0 & 0 \\
2.77 & 0 & 0 & 0 \\
0 & 0 & 0.69 & 2.08 \\
0 & 0 & 0 & 0
\end{pmatrix}
\]
\vfill
\begin{itemize}
\item Row 4 (\code{data}) is annihilated because it carries zero discriminative information across this corpus, whereas the rare term \code{eigenvalue} in Row 2 has its relative weight doubled!
\end{itemize}
\end{frame}


\begin{frame}{Example -- Column normalisation $D_{\text{IDF}} X D_{\text{len}}^{-1}$ (Part 3)}
\textbf{Step 3: Right-multiply $X_{\text{TF-IDF}}$ by $D_{\text{len}}^{-1}$ to normalise the columns in $\ell_2$ norm}:
\begin{itemize}
\item The Euclidean norms of the four columns of $X_{\text{TF-IDF}} = D_{\text{IDF}} X$ are:
\begin{align*}
\|\bx_1^{\text{TF-IDF}}\|_2 &= (\log 2)\sqrt{3^2 + 4^2} = 5\log 2, &
\|\bx_2^{\text{TF-IDF}}\|_2 &= (\log 2)\sqrt{2^2 + 0^2} = 2\log 2 \\
\|\bx_3^{\text{TF-IDF}}\|_2 &= (\log 2)\sqrt{1^2} = \log 2, &
\|\bx_4^{\text{TF-IDF}}\|_2 &= (\log 2)\sqrt{3^2} = 3\log 2
\end{align*}
\item Right-multiplying by $D_{\text{len}}^{-1} = (\log 2)^{-1}\diag(1/5,\; 1/2,\; 1,\; 1/3)$ cancels $\log 2$ and yields unit-norm columns:
\[
\tilde{X} = D_{\text{IDF}} X D_{\text{len}}^{-1} =
\begin{pmatrix}
3/5 & 1 & 0 & 0 \\
4/5 & 0 & 0 & 0 \\
0 & 0 & 1 & 1 \\
0 & 0 & 0 & 0
\end{pmatrix}
=
\begin{pmatrix}
0.6 & 1 & 0 & 0 \\
0.8 & 0 & 0 & 0 \\
0 & 0 & 1 & 1 \\
0 & 0 & 0 & 0
\end{pmatrix}
\]
\end{itemize}
\end{frame}


%%%%%%%%%%%%%%%%%%%%%
%%%%%%%%%%%%%%%%%%%%%
%%%%%%%%%%%%%%%%%%%%%
%%%%%%%%%%%%%%%%%%%%%
\SSsection{Angles and cosine similarity in high dimensions}

\begin{frame}{Why Euclidean distance fails for sparse document vectors}
Suppose we want to measure how similar two documents $\bx, \by \in \IRplus^m$ are. Why not use the Euclidean distance $\|\bx - \by\|_2$?
\vfill
\begin{itemize}
\item Consider three documents over $V = \{\code{matrix}, \code{vector}, \code{graph}\}$:
\[
\bx_1 = \begin{pmatrix} 1 \\ 1 \\ 0 \end{pmatrix}, \qquad
\bx_2 = \begin{pmatrix} 5 \\ 5 \\ 0 \end{pmatrix} = 5\bx_1, \qquad
\bx_3 = \begin{pmatrix} 0 \\ 0 \\ 1 \end{pmatrix}
\]
\item Documents $d_1$ and $d_2$ have the \textbf{exact same word proportions} (pure linear algebra), whereas $d_3$ shares \textbf{no words at all} with $d_1$
\vfill
\item Yet their Euclidean distances are:
\[
\|\bx_1 - \bx_2\|_2 = \left\|\begin{pmatrix} -4 \\ -4 \\ 0 \end{pmatrix}\right\|_2 = \sqrt{32} \approx 5.66
\quad > \quad
\|\bx_1 - \bx_3\|_2 = \left\|\begin{pmatrix} 1 \\ 1 \\ -1 \end{pmatrix}\right\|_2 = \sqrt{3} \approx 1.73
\]
\item Euclidean distance judges $d_1$ to be much closer to the unrelated short document $d_3$ than to a longer copy of itself!
\end{itemize}
\end{frame}


\begin{frame}{Direction matters more than magnitude}
\begin{center}
\begin{tikzpicture}[scale=0.85]
    \draw[->, thick] (0,0) -- (6.5,0) node[right, font=\small\bfseries] {Term 1 (\code{matrix})};
    \draw[->, thick] (0,0) -- (0,4.5) node[above, font=\small\bfseries] {Term 2 (\code{graph})};

    % Vector x1 and x2 = 3 x1
    \draw[->, ultra thick, blue!75!black] (0,0) -- (1.8,0.6) node[below right, font=\small] {$\bx_1$};
    \draw[->, thick, blue!75!black, dashed] (1.8,0.6) -- (5.4,1.8) node[right, font=\small] {$\bx_2 = 3\bx_1$};

    % Vector x3
    \draw[->, ultra thick, red!75!black] (0,0) -- (0.5,2.2) node[above left, font=\small] {$\bx_3$};

    % Angle arc
    \draw[thick, orange!90!black, ->] (1.1,0.367) arc (18.43:77.2:1.16);
    \node[orange!90!black, font=\small\bfseries] at (1.1,1.1) {$\theta$};

    % Euclidean distance lines
    \draw[dotted, thick, gray!80!black] (1.8,0.6) -- (0.5,2.2);
\end{tikzpicture}
\end{center}
\vfill
\begin{itemize}
\item Scaling a document vector by a positive constant ($\bx_2 = \alpha \bx_1$ with $\alpha > 0$) changes its length $\|\bx\|_2$ but preserves its \textbf{direction}
\item Therefore, semantic similarity between documents should be measured by the \defword{angle} $\theta$ between their vectors
\end{itemize}
\end{frame}


\begin{frame}{Cosine similarity}
\begin{importantdefinition}[Cosine similarity]
Let $\bx, \by \in \IR^m$ be two non-zero vectors. The \defword{cosine similarity} between $\bx$ and $\by$ is the cosine of the angle $\theta \in [0, \pi]$ between them:
\[
\mathsf{sim}(\bx, \by) = \cos(\theta) = \frac{\bx \bullet \by}{\|\bx\|_2 \, \|\by\|_2}
= \frac{\sum_{i=1}^m x_i y_i}{\sqrt{\sum_{i=1}^m x_i^2}\sqrt{\sum_{i=1}^m y_i^2}}
\]
Equivalently, in terms of the $\ell_2$-normalised unit vectors $\hat{\bx} = \frac{\bx}{\|\bx\|_2}$ and $\hat{\by} = \frac{\by}{\|\by\|_2}$,
\[
\mathsf{sim}(\bx, \by) = \hat{\bx} \bullet \hat{\by} = \hat{\bx}^T \hat{\by}
\]
\end{importantdefinition}
\end{frame}


\begin{frame}{Fundamental properties of cosine similarity}
\begin{proposition}[Properties of cosine similarity]\label{prop:cosine_properties}
Let $\bx, \by \in \IR^m$ be non-zero vectors. Then:
\begin{enumerate}
\item \textbf{Positive scale invariance}: For any scalars $\alpha > 0$ and $\beta > 0$,
\[
\mathsf{sim}(\alpha \bx, \beta \by) = \mathsf{sim}(\bx, \by)
\]
\item \textbf{General bounds}: By the Cauchy-Schwarz inequality, $\mathsf{sim}(\bx, \by) \in [-1, 1]$, with $\mathsf{sim}(\bx, \by) = 1 \iff \by = c\bx$ for some $c > 0$
\item \textbf{Non-negative orthant bounds}: If $\bx, \by \in \IRplus^m$ are non-negative vectors, then
\[
0 \le \mathsf{sim}(\bx, \by) \le 1
\]
Furthermore, for $\bx, \by \in \IRplus^m$,
\[
\mathsf{sim}(\bx, \by) = 0 \iff \mathsf{supp}(\bx) \cap \mathsf{supp}(\by) = \emptyset
\]
\end{enumerate}
\end{proposition}
\end{frame}


\begin{frame}[red]{Proof of Proposition~\ref{prop:cosine_properties} (Part 1: scale invariance)}
Let $\bx, \by \in \IR^m \setminus \{\b0\}$ and let $\alpha, \beta > 0$. Using the bilinearity of the dot product and the positive homogeneity of the Euclidean norm ($\|\alpha \bx\|_2 = |\alpha|\,\|\bx\|_2 = \alpha \|\bx\|_2$ since $\alpha > 0$), we compute
\begin{align*}
\mathsf{sim}(\alpha \bx, \beta \by)
&= \frac{(\alpha \bx) \bullet (\beta \by)}{\|\alpha \bx\|_2 \, \|\beta \by\|_2} \\
&= \frac{\alpha \beta (\bx \bullet \by)}{(\alpha \|\bx\|_2)(\beta \|\by\|_2)} \\
&= \frac{\cancel{\alpha \beta}(\bx \bullet \by)}{\cancel{\alpha \beta}\|\bx\|_2 \, \|\by\|_2} \\
&= \mathsf{sim}(\bx, \by)
\end{align*}
\vfill
(Note: if $\alpha < 0$ and $\beta > 0$, then $|\alpha| = -\alpha$, so $\mathsf{sim}(\alpha\bx, \beta\by) = -\mathsf{sim}(\bx, \by)$)\hfill\qed
\end{frame}


\begin{frame}[red]{Proof of Proposition~\ref{prop:cosine_properties} (Part 2: non-negative vectors)}
Assume $\bx, \by \in \IRplus^m \setminus \{\b0\}$, so $x_i \ge 0$ and $y_i \ge 0$ for every $i \in \{1, \ldots, m\}$
\vfill
\begin{enumerate}
\item Since the product of two non-negative real numbers is non-negative, $x_i y_i \ge 0$ for all $i = 1, \ldots, m$. Summing over $i$ yields
\[
\bx \bullet \by = \sum_{i=1}^m x_i y_i \ge 0
\]
Dividing by $\|\bx\|_2 \|\by\|_2 > 0$ and combining with Cauchy-Schwarz gives $0 \le \mathsf{sim}(\bx, \by) \le 1$
\vfill
\item Because each summand $x_i y_i$ is non-negative, a sum of non-negative terms $\sum_{i=1}^m x_i y_i$ equals $0$ if and only if \textbf{every individual term is zero}:
\[
\mathsf{sim}(\bx, \by) = 0 \iff \sum_{i=1}^m x_i y_i = 0 \iff \forall i \in \{1,\ldots,m\},\; x_i y_i = 0
\]
Since $x_i y_i = 0 \iff (x_i = 0 \textrm{ or } y_i = 0)$, there is no index $i$ with both $x_i > 0$ and $y_i > 0$, which means $\mathsf{supp}(\bx) \cap \mathsf{supp}(\by) = \emptyset$ \hfill\qed
\end{enumerate}
\end{frame}


\begin{frame}{Interpretation of Proposition~\ref{prop:cosine_properties} for text}
\begin{itemize}
\item Because bag-of-words and TF-IDF vectors lie in the non-negative orthant $\IRplus^m$, two documents can \textbf{never} have negative cosine similarity ($\theta \in [0, \pi/2]$)
\vfill
\item $\mathsf{sim}(\bx_i, \bx_j) = 0$ ($\theta = \pi/2$, orthogonal vectors) occurs \textbf{if and only if the two documents share no vocabulary words in common} (after any zero-IDF stop words are removed)
\vfill
\item Moreover, for $\ell_2$-normalised vectors $\hat{\bx} = \bx/\|\bx\|_2$ and $\hat{\by} = \by/\|\by\|_2$ on the unit sphere $S^{m-1}$, Euclidean distance and cosine similarity are monotonically related:
\[
\|\hat{\bx} - \hat{\by}\|_2^2 = \|\hat{\bx}\|_2^2 + \|\hat{\by}\|_2^2 - 2(\hat{\bx} \bullet \hat{\by}) = 2\left(1 - \mathsf{sim}(\bx, \by)\right)
\]
\end{itemize}
\end{frame}


\begin{frame}{All-pairs cosine similarity via the Gram matrix}
How do we compute the cosine similarity between \textbf{all pairs} of $n$ documents in a matrix $X = [\bx_1 \cdots \bx_n] \in \M_{m,n}$ at once?
\vfill
\begin{importanttheorem}[Cosine similarity Gram matrix]\label{th:cosine_gram}
Let $X = [\bx_1 \cdots \bx_n] \in \M_{m,n}$ have non-zero columns and let
\[
\hat{X} = X D_{\text{norm}}^{-1} = \begin{bmatrix} \hat{\bx}_1 & \hat{\bx}_2 & \cdots & \hat{\bx}_n \end{bmatrix} \in \M_{m,n}
\quad\textrm{where}\quad
D_{\text{norm}}^{-1} = \diag\!\left(\frac{1}{\|\bx_1\|_2}, \ldots, \frac{1}{\|\bx_n\|_2}\right)
\]
Then the symmetric \defword{Gram matrix}
\[
S = \hat{X}^T \hat{X} \in \M_n(\IR)
\]
has $(i,j)$ entry equal to the pairwise cosine similarity $s_{ij} = \hat{\bx}_i^T \hat{\bx}_j = \mathsf{sim}(\bx_i, \bx_j)$, with $s_{ii} = 1$ along the diagonal
\end{importanttheorem}
\end{frame}


\begin{frame}{Example -- Pairwise cosine similarity by hand (Part 1)}
Consider $n=3$ sparse document vectors in $\IRplus^5$:
\[
\bx_1 = \begin{pmatrix} 3 \\ 4 \\ 0 \\ 0 \\ 0 \end{pmatrix}, \qquad
\bx_2 = \begin{pmatrix} 6 \\ 8 \\ 0 \\ 0 \\ 0 \end{pmatrix}, \qquad
\bx_3 = \begin{pmatrix} 0 \\ 2 \\ 0 \\ 2 \\ 1 \end{pmatrix}
\]
\vfill
\textbf{Step 1: Compute the Euclidean norm of each document vector}:
\begin{align*}
\|\bx_1\|_2 &= \sqrt{3^2 + 4^2 + 0 + 0 + 0} = \sqrt{25} = 5 \\
\|\bx_2\|_2 &= \sqrt{6^2 + 8^2 + 0 + 0 + 0} = \sqrt{100} = 10 \\
\|\bx_3\|_2 &= \sqrt{0 + 2^2 + 0 + 2^2 + 1^2} = \sqrt{9} = 3
\end{align*}
\end{frame}


\begin{frame}{Example -- Pairwise cosine similarity by hand (Part 2)}
\textbf{Step 2: Compute pairwise dot products (only shared non-zero indices contribute!)}:
\begin{align*}
\bx_1 \bullet \bx_2 &= 3(6) + 4(8) = 18 + 32 = 50 \\
\bx_1 \bullet \bx_3 &= 4(2) = 8 \\
\bx_2 \bullet \bx_3 &= 8(2) = 16
\end{align*}
\vfill
\textbf{Step 3: Divide by norms to form the Gram matrix $S = \hat{X}^T \hat{X}$}:
\[
\mathsf{sim}(\bx_1, \bx_2) = \frac{50}{5 \times 10} = 1, \quad
\mathsf{sim}(\bx_1, \bx_3) = \frac{8}{5 \times 3} = \frac{8}{15}, \quad
\mathsf{sim}(\bx_2, \bx_3) = \frac{16}{10 \times 3} = \frac{8}{15}
\]
\[
\implies S = \hat{X}^T \hat{X} =
\begin{pmatrix}
1 & 1 & 8/15 \\
1 & 1 & 8/15 \\
8/15 & 8/15 & 1
\end{pmatrix}
\approx
\begin{pmatrix}
1.000 & 1.000 & 0.533 \\
1.000 & 1.000 & 0.533 \\
0.533 & 0.533 & 1.000
\end{pmatrix}
\]
\end{frame}


%%%%%%%%%%%%%%%%%%%%%
%%%%%%%%%%%%%%%%%%%%%
%%%%%%%%%%%%%%%%%%%%%
%%%%%%%%%%%%%%%%%%%%%
\SSsection{Non-negative matrix factorisation (NMF)}

%%%%%%%%%%%%%%%%%%%%%
%%%%%%%%%%%%%%%%%%%%%
\SSsubsection{SVD vs NMF on non-negative data}

\begin{frame}{Recall from Unit 04: low-rank approximation via SVD}
In Unit 04, we saw that any matrix $X \in \M_{m,n}(\IR)$ admits a Singular Value Decomposition $X = U \Sigma V^T = \sum_{r=1}^{\rank(X)} \sigma_r \bu_r \bv_r^T$
\vfill
\begin{itemize}
\item By the \textbf{Eckart-Young theorem}, the best rank-$k$ approximation to $X$ in the Frobenius norm $\|A\|_F = \sqrt{\sum_{i=1}^m \sum_{j=1}^n a_{ij}^2}$ is the \defword{truncated SVD}:
\[
X_k = \sum_{r=1}^k \sigma_r \bu_r \bv_r^T = U_k \Sigma_k V_k^T
\]
\vfill
\item Applied to a term-document matrix $X$, truncated SVD is known in information retrieval as \defword{Latent Semantic Analysis} (\defword{LSA}, Deerwester et al.\ 1990)
\end{itemize}
\end{frame}


\begin{frame}{Why truncated SVD is awkward on non-negative data}
Many real-world data matrices are inherently \textbf{non-negative} ($x_{ij} \ge 0$):
\begin{itemize}
\item Term-document word count and TF-IDF matrices
\item Pixel intensities in greyscale or colour images
\item Audio spectrograms, chemical concentrations and gene expression counts
\end{itemize}
\vfill
Why does truncated SVD struggle with interpretability on $X \in \M_{m,n}(\IRplus)$?
\begin{itemize}
\item By the Perron-Frobenius theorem, the first singular vectors $\bu_1, \bv_1$ can be chosen non-negative
\item However, $U_k$ and $V_k$ have \textbf{orthonormal columns}: for every $r \ge 2$,
\[
\bu_r \bullet \bu_1 = 0 \quad\textrm{and}\quad \bv_r \bullet \bv_1 = 0
\]
\item As proved in Proposition~\ref{prop:cosine_properties}, a non-zero vector orthogonal to a strictly positive vector $\bu_1$ \textbf{must contain both positive and negative entries}!
\end{itemize}
\end{frame}


\begin{frame}{Subtractive cancellation vs additive parts}
\begin{itemize}
\item Because $\bu_2, \ldots, \bu_k$ and $\bv_2, \ldots, \bv_k$ contain negative numbers, SVD reconstructs a document $\bx_j \approx \sum_{r=1}^k (\sigma_r v_{jr}) \bu_r$ through \textbf{subtractive cancellation}:
\begin{itemize}
\item What does it mean for a ``topic'' vector $\bu_2$ to contain $-0.42 \times \code{"matrix"}$?
\item What does it mean to subtract $0.3$ of Topic 2 from a document so that negative word counts cancel out?
\end{itemize}
\vfill
\item For example, on our $6 \times 4$ corpus matrix $X$, the first two left singular vectors in $U_2$ are
\[
\bu_1 = \begin{pmatrix} 0.30 \\ 0.30 \\ 0.11 \\ 0.32 \\ 0.31 \\ 0.78 \end{pmatrix} \ge \b0, \qquad
\bu_2 = \begin{pmatrix} \boldred{-0.54} \\ \boldred{-0.47} \\ \boldred{-0.24} \\ 0.45 \\ 0.45 \\ 0.14 \end{pmatrix} \notin \IRplus^6
\]
\item To obtain directly interpretable topics (or parts of an image), we drop the orthogonality constraint and instead enforce \defword{non-negativity} on both factors!
\end{itemize}
\end{frame}


\begin{frame}{Non-negative Matrix Factorisation (NMF)}
\begin{importantdefinition}[Non-negative Matrix Factorisation (NMF)]
Given a non-negative matrix $X \in \M_{m,n}(\IRplus)$ and a target rank $k < \min(m, n)$, \defword{Non-negative Matrix Factorisation} (\defword{NMF}, Paatero \& Tapper 1994; Lee \& Seung 1999) seeks two non-negative matrices
\[
W \in \M_{m,k}(\IRplus) \quad\textrm{and}\quad H \in \M_{k,n}(\IRplus)
\]
such that $X \approx W H$, by solving the constrained optimisation problem
\[
\min_{W \in \M_{m,k}(\IRplus),\; H \in \M_{k,n}(\IRplus)} \|X - W H\|_F^2
= \sum_{i=1}^m \sum_{j=1}^n \left(x_{ij} - \sum_{r=1}^k w_{ir} h_{rj}\right)^2
\]
subject to $w_{ir} \ge 0$ and $h_{rj} \ge 0$ for all $i, r, j$
\end{importantdefinition}
\end{frame}


\begin{frame}{Visualising the low-rank NMF decomposition $X \approx W H$}
\begin{center}
\begin{tikzpicture}[scale=0.9, every node/.style={font=\small}]
    % X matrix
    \draw[thick, fill=blue!15] (0,0) rectangle (3.0,3.6);
    \node[font=\large\bfseries] at (1.5,1.8) {$X \ge 0$};
    \node[left] at (0,1.8) {$m$ terms};
    \node[above] at (1.5,3.6) {$n$ documents};
    \node[below] at (1.5,0) {$m \times n$};

    \node[font=\Large\bfseries] at (3.7,1.8) {$\approx$};

    % W matrix
    \draw[thick, fill=orange!20] (4.4,0) rectangle (5.6,3.6);
    \node[font=\large\bfseries] at (5.0,1.8) {$W \ge 0$};
    \node[above] at (5.0,3.6) {$k$ topics};
    \node[below] at (5.0,0) {$m \times k$};

    % H matrix
    \draw[thick, fill=green!20] (6.1,2.4) rectangle (9.1,3.6);
    \node[font=\large\bfseries] at (7.6,3.0) {$H \ge 0$};
    \node[above] at (7.6,3.6) {$n$ documents};
    \node[right] at (9.1,3.0) {$k$ topics};
    \node[below] at (7.6,2.4) {$k \times n$};
\end{tikzpicture}
\end{center}
\vfill
\begin{itemize}
\item When $k \ll \min(m, n)$, NMF compresses $mn$ entries into $k(m + n)$ non-negative parameters while preserving an \textbf{additive, parts-based} structure
\end{itemize}
\end{frame}


\begin{frame}{Comparison: truncated SVD vs NMF}
\begin{center}
\begin{tabular}{lll}
\toprule
\textbf{Property} & \textbf{Truncated SVD ($U_k \Sigma_k V_k^T$)} & \textbf{NMF ($W H$)} \\
\midrule
Structural constraint & Orthogonality ($U_k^T U_k = V_k^T V_k = \II_k$) & Non-negativity ($W \ge 0, H \ge 0$) \\
Entries of factors & Positive and negative ($\in \IR$) & Non-negative only ($\in \IRplus$) \\
Reconstruction & Subtractive cancellation & Purely additive parts \\
Frobenius error $\|X - \hat{X}\|_F$ & Globally minimal for rank $k$ & Higher or equal ($\ge \|X - X_k\|_F$) \\
Optimisation problem & Exact via eigendecomposition & Non-convex; iterative local minima \\
Nested factors across $k$ & Yes (first $k$ cols stay fixed) & No (recomputed for each $k$) \\
\bottomrule
\end{tabular}
\end{center}
\end{frame}


%%%%%%%%%%%%%%%%%%%%%
%%%%%%%%%%%%%%%%%%%%%
\SSsubsection{Topic modelling interpretation of W and H}

\begin{frame}{Column-wise view: documents as additive mixtures of topics}
Write $W = \begin{bmatrix} \bw_1 & \cdots & \bw_k \end{bmatrix} \in \M_{m,k}(\IRplus)$ in terms of its columns $\bw_r \in \IRplus^m$ and $H = \begin{bmatrix} \bh_1 & \cdots & \bh_n \end{bmatrix} \in \M_{k,n}(\IRplus)$ in terms of its columns $\bh_j \in \IRplus^k$
\vfill
By standard column-wise matrix multiplication, the $j$-th column of $X \approx W H$ is
\[
\bx_j \approx W \bh_j = h_{1j} \bw_1 + h_{2j} \bw_2 + \cdots + h_{kj} \bw_k = \sum_{r=1}^k h_{rj} \bw_r
\]
\vfill
\begin{itemize}
\item \textbf{Columns of $W$ ($\bw_r \in \IRplus^m$, the dictionary / topic basis)}:
\begin{itemize}
\item Column $\bw_r$ gives the weight of each of the $m$ vocabulary words in \defword{topic $r$}
\item The largest entries of $\bw_r$ identify the signature keywords of topic $r$
\end{itemize}
\vfill
\item \textbf{Columns of $H$ ($\bh_j \in \IRplus^k$, the document topic encodings)}:
\begin{itemize}
\item Entry $h_{rj} \ge 0$ tells how strongly document $d_j$ expresses topic $r$
\end{itemize}
\end{itemize}
\end{frame}


\begin{frame}{Geometric interpretation: the conical hull of $W$}
\begin{columns}[c]
\begin{column}{0.48\textwidth}
\begin{center}
\begin{tikzpicture}[scale=0.85]
    % Axes in R^3
    \draw[->, thick] (0,0) -- (4.2,0) node[right, font=\small] {$w_1$};
    \draw[->, thick] (0,0) -- (0,3.8) node[above, font=\small] {$w_3$};
    \draw[->, thick] (0,0) -- (-1.8,-1.4) node[below left, font=\small] {$w_2$};

    % Shaded cone between w_1 and w_2
    \fill[orange!25, opacity=0.7] (0,0) -- (3.4,0.8) -- (0.8,3.2) -- cycle;

    % Topic vectors
    \draw[->, ultra thick, orange!85!black] (0,0) -- (3.4,0.8) node[right, font=\small\bfseries] {$\bw_1$ (Topic 1)};
    \draw[->, ultra thick, orange!85!black] (0,0) -- (0.8,3.2) node[above right, font=\small\bfseries] {$\bw_2$ (Topic 2)};

    % Document points inside the cone
    \fill[blue!75!black] (2.4,1.1) circle (2.5pt) node[below, font=\footnotesize] {$W\bh_1$};
    \fill[blue!75!black] (1.3,2.3) circle (2.5pt) node[left, font=\footnotesize] {$W\bh_2$};
    \fill[blue!75!black] (2.0,1.8) circle (2.5pt) node[right, font=\footnotesize] {$W\bh_3$};
\end{tikzpicture}
\end{center}
\end{column}
\begin{column}{0.52\textwidth}
\begin{itemize}
\item In unconstrained linear algebra, $\{W\bh : \bh \in \IR^k\}$ is the $k$-dimensional subspace $\Span(\bw_1, \ldots, \bw_k)$
\vfill
\item In NMF, because $\bh_j \in \IRplus^k$ has non-negative weights $h_{rj} \ge 0$, the reconstructed documents $W\bh_j$ must lie inside the \defword{conical hull} (or simplicial cone) generated by the columns of $W$:
\[
\mathsf{cone}(W) = \left\{\sum_{r=1}^k h_r \bw_r : h_1, \ldots, h_k \ge 0\right\} \subset \IRplus^m
\]
\end{itemize}
\end{column}
\end{columns}
\end{frame}


\begin{frame}{Rank-one outer product expansion and polysemy}
Just as the SVD writes $X_k = \sum_{r=1}^k \sigma_r \bu_r \bv_r^T$ as a sum of $k$ rank-one matrices, NMF decomposes $W H$ into $k$ \textbf{non-negative rank-one layers}:
\[
W H = \sum_{r=1}^k \bw_r \tilde{\bh}_r^T
\]
where $\bw_r \in \IRplus^m$ is the $r$-th column of $W$ and $\tilde{\bh}_r^T \in \IRplus^{1 \times n}$ is the $r$-th row of $H$
\vfill
\begin{itemize}
\item Each outer product $M_r = \bw_r \tilde{\bh}_r^T \in \M_{m,n}(\IRplus)$ displays the contribution of \textbf{topic $r$ alone} across all $m$ words and all $n$ documents
\item Because columns of $W$ do not have to be orthogonal, a single word $w_i$ with multiple meanings (\defword{polysemy}, e.g., \code{"degree"} in polynomial algebra vs vertex degree in graphs) can have positive weight $w_{ir} > 0$ in \textbf{more than one topic column}!
\end{itemize}
\end{frame}


\begin{frame}{Example -- Interpreting $W$ and $H$ (Part 1)}
Suppose a corpus of $n=4$ documents over $m=5$ terms $\{\code{matrix}, \code{vector}, \code{rank}, \code{graph}, \code{vertex}\}$ has term-document matrix $X \in \M_{5,4}(\IRplus)$ and rank-$2$ NMF gives $X \approx W H$:
\[
X = \begin{pmatrix}
6 & 4 & 0 & 2 \\
3 & 2 & 0 & 1 \\
3 & 2 & 0 & 1 \\
0 & 0 & 6 & 4 \\
0 & 0 & 3 & 2
\end{pmatrix},
\qquad
W = \begin{pmatrix}
2 & 0 \\
1 & 0 \\
1 & 0 \\
0 & 2 \\
0 & 1
\end{pmatrix},
\qquad
H = \begin{pmatrix}
3 & 2 & 0 & 1 \\
0 & 0 & 3 & 2
\end{pmatrix}
\]
(where rows $1,\ldots,5$ of $X$ and $W$ are $\code{matrix}, \code{vector}, \code{rank}, \code{graph}, \code{vertex}$ while columns $1,\ldots,4$ of $X$ and $H$ are $d_1, d_2, d_3, d_4$)
\vfill
\textbf{Questions}:
\begin{enumerate}
\item How should we label Topic 1 ($\bw_1$) and Topic 2 ($\bw_2$) from $W$?
\item How should we characterise documents $d_1, d_2, d_3, d_4$ from $H$?
\item What is the reconstruction error $\|X - WH\|_F^2$?
\end{enumerate}
\end{frame}


\begin{frame}{Example -- Interpreting $W$ and $H$ (Part 2)}
\begin{enumerate}
\item \textbf{Labelling topics from the columns of $W$}:
\begin{itemize}
\item Column $\bw_1 = (2, 1, 1, 0, 0)^T$ has non-zero weights only on \code{matrix}, \code{vector} and \code{rank} $\implies$ \textbf{Topic 1 is Linear Algebra}
\item Column $\bw_2 = (0, 0, 0, 2, 1)^T$ has non-zero weights only on \code{graph} and \code{vertex} $\implies$ \textbf{Topic 2 is Graph Theory}
\end{itemize}
\vfill
\item \textbf{Classifying documents from the columns of $H$}:
\begin{itemize}
\item $\bh_1 = (3, 0)^T$ and $\bh_2 = (2, 0)^T$ load purely on Topic 1 (Linear Algebra)
\item $\bh_3 = (0, 3)^T$ loads purely on Topic 2 (Graph Theory)
\item $\bh_4 = (1, 2)^T$ is a \textbf{multi-topic document}: $\bx_4 = 1\bw_1 + 2\bw_2$ ($1/3$ Linear Algebra, $2/3$ Graph Theory, e.g., spectral graph theory!)
\end{itemize}
\vfill
\item \textbf{Reconstruction error}: Multiplying $W H$ reproduces $X$ exactly, so $\|X - WH\|_F^2 = 0$ and $\rank(X) = 2$
\end{enumerate}
\end{frame}


%%%%%%%%%%%%%%%%%%%%%
%%%%%%%%%%%%%%%%%%%%%
\SSsubsection{Alternating least squares (ALS)}

\begin{frame}{Why computing NMF is an optimisation challenge}
Consider the NMF objective function
\[
f(W, H) = \|X - W H\|_F^2 = \sum_{i=1}^m \sum_{j=1}^n \left(x_{ij} - \sum_{r=1}^k w_{ir} h_{rj}\right)^2
\]
\vfill
\begin{enumerate}
\item \textbf{Non-uniqueness (scaling and permutation ambiguity)}:
\begin{itemize}
\item For any permutation matrix $P \in \M_k$ and any positive diagonal matrix $D = \diag(d_1, \ldots, d_k)$ with $d_r > 0$, setting $\tilde{W} = W P D \ge 0$ and $\tilde{H} = D^{-1} P^T H \ge 0$ gives $\tilde{W}\tilde{H} = W H$
\end{itemize}
\vfill
\item \textbf{Non-convexity jointly in $(W, H)$}:
\begin{itemize}
\item Because unknown parameters $w_{ir}$ and $h_{rj}$ are multiplied together in $W H$, $f(W, H)$ is \textbf{not convex} in $(W, H)$ simultaneously and can have multiple local minima
\end{itemize}
\end{enumerate}
\end{frame}


\begin{frame}{The key linear algebra insight: bi-convexity}
Although $f(W, H) = \|X - W H\|_F^2$ is non-convex jointly in $(W, H)$, it is \defword{bi-convex} (convex in $H$ when $W$ is fixed and convex in $W$ when $H$ is fixed)!
\vfill
\begin{enumerate}
\item \textbf{Hold $W \in \M_{m,k}(\IRplus)$ fixed}:
\[
\|X - W H\|_F^2 = \sum_{j=1}^n \|\bx_j - W \bh_j\|_2^2
\]
Each term $\|\bx_j - W \bh_j\|_2^2$ is a standard \textbf{linear least-squares problem} $W \bh_j \approx \bx_j$ of the exact form studied in Unit 03 and Unit 04!
\vfill
\item \textbf{Hold $H \in \M_{k,n}(\IRplus)$ fixed}: Using $\|A\|_F^2 = \|A^T\|_F^2$ and $(WH)^T = H^T W^T$,
\[
\|X - W H\|_F^2 = \|X^T - H^T W^T\|_F^2
\]
This is again a linear least-squares problem for the columns of $W^T$ with design matrix $H^T \in \M_{n,k}$!
\end{enumerate}
\end{frame}


\begin{frame}{Deriving the projected least-squares update for $H$}
\begin{proposition}[Projected least-squares step for $H$]\label{prop:als_H_step}
Suppose $W \in \M_{m,k}(\IRplus)$ is fixed and has full column rank ($\rank(W) = k$, so $W^T W \in \M_k$ is invertible). Then:
\begin{enumerate}
\item For each column $j \in \{1, \ldots, n\}$, the unconstrained least-squares problem $\min_{\bh_j \in \IR^k} \|\bx_j - W \bh_j\|_2^2$ has the unique normal-equations solution
\[
\bh_j^{\text{LS}} = (W^T W)^{-1} W^T \bx_j
\]
\item Stacking all $n$ columns side by side into $H_{\text{LS}} = [\bh_1^{\text{LS}} \cdots \bh_n^{\text{LS}}] \in \M_{k,n}(\IR)$ gives the matrix formula
\[
H_{\text{LS}} = (W^T W)^{-1} W^T X
\]
\item In projected \defword{Alternating Least Squares} (\defword{ALS}), we enforce non-negativity by clipping negative entries to zero entrywise:
\[
H \leftarrow \max(0, H_{\text{LS}}) \quad\textrm{i.e.,}\quad h_{rj} = \max\!\left(0, (H_{\text{LS}})_{rj}\right)
\]
\end{enumerate}
\end{proposition}
\end{frame}


\begin{frame}{Deriving the projected least-squares update for $W$}
\begin{proposition}[Projected least-squares step for $W$]\label{prop:als_W_step}
Suppose $H \in \M_{k,n}(\IRplus)$ is fixed and has full row rank ($\rank(H) = k$, so $H H^T \in \M_k$ is invertible). Then:
\begin{enumerate}
\item Minimising $\|X - W H\|_F^2 = \|X^T - H^T W^T\|_F^2$ without sign constraints yields the normal equations for $W^T$:
\[
(H^T)^T (H^T) W_{\text{LS}}^T = (H^T)^T X^T \iff (H H^T) W_{\text{LS}}^T = H X^T
\]
\item Solving for $W_{\text{LS}}^T$ and taking the transpose (using $((H H^T)^{-1})^T = (H H^T)^{-1}$) gives
\[
W_{\text{LS}}^T = (H H^T)^{-1} H X^T \iff W_{\text{LS}} = X H^T (H H^T)^{-1}
\]
\item Clipping negative entries to zero entrywise gives the projected ALS update:
\[
W \leftarrow \max(0, W_{\text{LS}})
\]
\end{enumerate}
\end{proposition}
\end{frame}


\begin{frame}{The Alternating Least Squares (ALS) algorithm for NMF}
Putting Propositions~\ref{prop:als_H_step} and~\ref{prop:als_W_step} together yields the \defword{Projected Alternating Least Squares} (\defword{ALS}) algorithm (Paatero \& Tapper 1994; Berry et al.\ 2007):
\vfill
\begin{enumerate}
\item \textbf{Input}: Non-negative matrix $X \in \M_{m,n}(\IRplus)$ and target rank $k < \min(m, n)$
\item \textbf{Initialise}: Choose a non-negative initial matrix $W^{(0)} \in \M_{m,k}(\IRplus)$ with $\rank(W^{(0)}) = k$ (e.g., using positive random numbers)
\item \textbf{Repeat for $t = 0, 1, 2, \ldots$ until convergence}:
\begin{enumerate}
\item \textbf{Update $H$ (fixing $W^{(t)}$)}: Solve $(W^{(t)T} W^{(t)}) H_{\text{LS}} = W^{(t)T} X$ and project:
\[
H^{(t+1)} = \max\!\left(0,\; (W^{(t)T} W^{(t)})^{-1} W^{(t)T} X\right)
\]
\item \textbf{Update $W$ (fixing $H^{(t+1)}$)}: Solve $W_{\text{LS}} (H^{(t+1)} H^{(t+1)T}) = X H^{(t+1)T}$ and project:
\[
W^{(t+1)} = \max\!\left(0,\; X H^{(t+1)T} (H^{(t+1)} H^{(t+1)T})^{-1}\right)
\]
\item \emph{(Optional)} Normalise each column of $W^{(t+1)}$ to unit $\ell_2$ or $\ell_1$ norm to prevent scaling drift
\end{enumerate}
\end{enumerate}
\end{frame}


\begin{frame}{Example -- One ALS update step by hand (Part 1)}
Let $X \in \M_{3,3}(\IRplus)$ and let the initial rank-$2$ dictionary matrix $W^{(0)} \in \M_{3,2}(\IRplus)$ be
\[
X = \begin{pmatrix} 3 & 0 & 2 \\ 0 & 3 & 3 \\ 0 & 3 & 1 \end{pmatrix}, \qquad
W^{(0)} = \begin{pmatrix} 1 & 0 \\ 1 & 1 \\ 0 & 1 \end{pmatrix}
\]
Let us compute the first ALS update $H^{(1)} = \max\!\left(0, (W^{(0)T} W^{(0)})^{-1} W^{(0)T} X\right)$ by hand
\vfill
\textbf{Step 1: Compute $W^{(0)T} W^{(0)} \in \M_2$ and its inverse}:
\[
W^{(0)T} W^{(0)} =
\begin{pmatrix} 1 & 1 & 0 \\ 0 & 1 & 1 \end{pmatrix}
\begin{pmatrix} 1 & 0 \\ 1 & 1 \\ 0 & 1 \end{pmatrix}
= \begin{pmatrix} 2 & 1 \\ 1 & 2 \end{pmatrix}
\]
Since $\det(W^{(0)T} W^{(0)}) = 2(2) - 1(1) = 3 \neq 0$, the $2 \times 2$ inverse formula gives
\[
(W^{(0)T} W^{(0)})^{-1} = \frac{1}{3} \begin{pmatrix} 2 & -1 \\ -1 & 2 \end{pmatrix}
\]
\end{frame}


\begin{frame}{Example -- One ALS update step by hand (Part 2)}
\textbf{Step 2: Compute $W^{(0)T} X \in \M_{2,3}$}:
\[
W^{(0)T} X =
\begin{pmatrix} 1 & 1 & 0 \\ 0 & 1 & 1 \end{pmatrix}
\begin{pmatrix} 3 & 0 & 2 \\ 0 & 3 & 3 \\ 0 & 3 & 1 \end{pmatrix}
= \begin{pmatrix} 3 & 3 & 5 \\ 0 & 6 & 4 \end{pmatrix}
\]
\vfill
\textbf{Step 3: Compute the unconstrained least-squares matrix $H_{\text{LS}}^{(1)}$}:
\[
H_{\text{LS}}^{(1)} = (W^{(0)T} W^{(0)})^{-1} W^{(0)T} X
= \frac{1}{3} \begin{pmatrix} 2 & -1 \\ -1 & 2 \end{pmatrix} \begin{pmatrix} 3 & 3 & 5 \\ 0 & 6 & 4 \end{pmatrix}
= \frac{1}{3} \begin{pmatrix} 6 & 0 & 6 \\ -3 & 9 & 3 \end{pmatrix}
= \begin{pmatrix} 2 & 0 & 2 \\ \boldred{-1} & 3 & 1 \end{pmatrix}
\]
\vfill
\textbf{Step 4: Project onto $\IRplus$ by clipping negative entries to zero}:
\[
H^{(1)} = \max\!\left(0, H_{\text{LS}}^{(1)}\right) = \begin{pmatrix} 2 & 0 & 2 \\ \boldred{0} & 3 & 1 \end{pmatrix} \in \M_{2,3}(\IRplus)
\]
Notice how unconstrained least squares produced $-1$ in entry $(2,1)$, which the projection step clips to $0$!
\end{frame}


\begin{frame}{Example -- Evaluating the reconstruction $W^{(0)} H^{(1)}$ (Part 3)}
\textbf{Step 5: Compute the rank-$2$ approximation $\hat{X}^{(1)} = W^{(0)} H^{(1)}$ and its Frobenius error}:
\[
\hat{X}^{(1)} = W^{(0)} H^{(1)} =
\begin{pmatrix} 1 & 0 \\ 1 & 1 \\ 0 & 1 \end{pmatrix}
\begin{pmatrix} 2 & 0 & 2 \\ 0 & 3 & 1 \end{pmatrix}
= \begin{pmatrix} 2 & 0 & 2 \\ 2 & 3 & 3 \\ 0 & 3 & 1 \end{pmatrix}
\]
\vfill
Comparing $\hat{X}^{(1)}$ with $X = \begin{pmatrix} 3 & 0 & 2 \\ 0 & 3 & 3 \\ 0 & 3 & 1 \end{pmatrix}$, the residual matrix is
\[
X - W^{(0)} H^{(1)} = \begin{pmatrix} 1 & 0 & 0 \\ -2 & 0 & 0 \\ 0 & 0 & 0 \end{pmatrix}
\implies
\|X - W^{(0)} H^{(1)}\|_F^2 = 1^2 + (-2)^2 = 5
\]
\vfill
\begin{itemize}
\item Columns 2 and 3 of $X$ are reconstructed \textbf{with zero error} on the very first half-step; the next half-step updates $W^{(1)} = \max\!\left(0, X H^{(1)T}(H^{(1)}H^{(1)T})^{-1}\right)$ to reduce the error on column 1 even further!
\end{itemize}
\end{frame}


\begin{frame}{Remark: exact NNLS and multiplicative updates}
\begin{itemize}
\item \textbf{Projected ALS vs exact Non-Negative Least Squares (NNLS)}:
\begin{itemize}
\item Solving unconstrained least squares and clipping $\max(0, \cdot)$ is fast (requiring only a $k \times k$ linear system $W^T W$ per iteration) and works well when $k$ is small and $X$ is sparse
\item A more accurate variant solves each column subproblem $\min_{\bh_j \ge \b0} \|\bx_j - W \bh_j\|_2^2$ using an exact active-set \defword{NNLS} solver (Lawson \& Hanson 1974), guaranteeing that $\|X - W^{(t)}H^{(t)}\|_F^2$ decreases monotonically at every step
\end{itemize}
\vfill
\item \textbf{Lee \& Seung (1999) multiplicative updates}:
\begin{itemize}
\item Another popular algorithm updates entries by elementwise multiplication:
\[
h_{rj} \leftarrow h_{rj} \frac{(W^T X)_{rj}}{(W^T W H)_{rj}}, \qquad
w_{ir} \leftarrow w_{ir} \frac{(X H^T)_{ir}}{(W H H^T)_{ir}}
\]
\item Because all factors on the right-hand side are non-negative, $W$ and $H$ automatically remain non-negative without clipping while $\|X - WH\|_F^2$ is guaranteed not to increase
\end{itemize}
\end{itemize}
\end{frame}


\begin{frame}[fragile]{R demonstration -- TF-IDF and cosine similarity}
Let us build an $8 \times 6$ term-document matrix in R with $6$ documents (3 on linear algebra, 3 on graphs) and compute $\tilde{X} = D_{\text{IDF}} X D_{\text{len}}^{-1}$ and $S = \hat{X}^T \hat{X}$:
\vfill
\begin{knitrout}
\definecolor{shadecolor}{rgb}{0.969, 0.969, 0.969}\color{fgcolor}\begin{kframe}
\begin{alltt}
\hldef{terms} \hlkwb{=} \hlkwd{c}\hldef{(}\hlsng{"matrix"}\hldef{,} \hlsng{"vector"}\hldef{,} \hlsng{"eigen"}\hldef{,} \hlsng{"rank"}\hldef{,} \hlsng{"graph"}\hldef{,} \hlsng{"vertex"}\hldef{,} \hlsng{"edge"}\hldef{,} \hlsng{"data"}\hldef{)}
\hldef{X} \hlkwb{=} \hlkwd{matrix}\hldef{(}\hlkwd{c}\hldef{(}
  \hlnum{4}\hldef{,} \hlnum{3}\hldef{,} \hlnum{2}\hldef{,} \hlnum{1}\hldef{,} \hlnum{0}\hldef{,} \hlnum{0}\hldef{,} \hlnum{0}\hldef{,} \hlnum{5}\hldef{,}  \hlcom{# d1: linear algebra}
  \hlnum{3}\hldef{,} \hlnum{4}\hldef{,} \hlnum{0}\hldef{,} \hlnum{2}\hldef{,} \hlnum{0}\hldef{,} \hlnum{0}\hldef{,} \hlnum{0}\hldef{,} \hlnum{4}\hldef{,}  \hlcom{# d2: linear algebra}
  \hlnum{2}\hldef{,} \hlnum{1}\hldef{,} \hlnum{3}\hldef{,} \hlnum{2}\hldef{,} \hlnum{0}\hldef{,} \hlnum{0}\hldef{,} \hlnum{0}\hldef{,} \hlnum{6}\hldef{,}  \hlcom{# d3: linear algebra}
  \hlnum{0}\hldef{,} \hlnum{0}\hldef{,} \hlnum{0}\hldef{,} \hlnum{0}\hldef{,} \hlnum{4}\hldef{,} \hlnum{3}\hldef{,} \hlnum{2}\hldef{,} \hlnum{5}\hldef{,}  \hlcom{# d4: graph theory}
  \hlnum{0}\hldef{,} \hlnum{0}\hldef{,} \hlnum{0}\hldef{,} \hlnum{0}\hldef{,} \hlnum{3}\hldef{,} \hlnum{2}\hldef{,} \hlnum{4}\hldef{,} \hlnum{4}\hldef{,}  \hlcom{# d5: graph theory}
  \hlnum{0}\hldef{,} \hlnum{1}\hldef{,} \hlnum{0}\hldef{,} \hlnum{0}\hldef{,} \hlnum{2}\hldef{,} \hlnum{4}\hldef{,} \hlnum{3}\hldef{,} \hlnum{5}   \hlcom{# d6: graph theory (mentions "vector" once)}
\hldef{),} \hlkwc{nrow} \hldef{=} \hlnum{8}\hldef{,} \hlkwc{ncol} \hldef{=} \hlnum{6}\hldef{,} \hlkwc{dimnames} \hldef{=} \hlkwd{list}\hldef{(terms,} \hlkwd{paste0}\hldef{(}\hlsng{"d"}\hldef{,} \hlnum{1}\hlopt{:}\hlnum{6}\hldef{)))}

\hlcom{# Diagonal IDF scaling: D_IDF %*% X}
\hldef{df_i} \hlkwb{=} \hlkwd{rowSums}\hldef{(X} \hlopt{>} \hlnum{0}\hldef{)}
\hldef{idf_i} \hlkwb{=} \hlkwd{log}\hldef{(}\hlkwd{ncol}\hldef{(X)} \hlopt{/} \hldef{df_i)}
\hldef{X_tfidf} \hlkwb{=} \hlkwd{diag}\hldef{(idf_i)} \hlopt{%*%} \hldef{X}

\hlcom{# Column l2 normalisation: X_hat = X_tfidf %*% D_norm^\{-1\}}
\hldef{col_norms} \hlkwb{=} \hlkwd{sqrt}\hldef{(}\hlkwd{colSums}\hldef{(X_tfidf}\hlopt{^}\hlnum{2}\hldef{))}
\hldef{X_hat} \hlkwb{=} \hldef{X_tfidf} \hlopt{%*%} \hlkwd{diag}\hldef{(}\hlnum{1} \hlopt{/} \hldef{col_norms)}
\hldef{S_gram} \hlkwb{=} \hlkwd{round}\hldef{(}\hlkwd{t}\hldef{(X_hat)} \hlopt{%*%} \hldef{X_hat,} \hlnum{2}\hldef{)}
\end{alltt}
\end{kframe}
\end{knitrout}
\end{frame}




\maxFrameImage{FIGS/U11-plot-cosine-gram-1.pdf}


\begin{frame}[fragile]{R demonstration -- Projected ALS for NMF}
We now run $30$ iterations of projected Alternating Least Squares with rank $k = 2$ on our TF-IDF matrix:
\vfill
\begin{knitrout}
\definecolor{shadecolor}{rgb}{0.969, 0.969, 0.969}\color{fgcolor}\begin{kframe}
\begin{alltt}
\hlkwd{set.seed}\hldef{(}\hlnum{2740}\hldef{)}
\hldef{k} \hlkwb{=} \hlnum{2}
\hldef{W} \hlkwb{=} \hlkwd{matrix}\hldef{(}\hlkwd{runif}\hldef{(}\hlkwd{nrow}\hldef{(X_tfidf)} \hlopt{*} \hldef{k,} \hlnum{0.1}\hldef{,} \hlnum{1}\hldef{),} \hlkwc{nrow} \hldef{=} \hlkwd{nrow}\hldef{(X_tfidf),} \hlkwc{ncol} \hldef{= k)}
\hldef{H} \hlkwb{=} \hlkwd{matrix}\hldef{(}\hlnum{0}\hldef{,} \hlkwc{nrow} \hldef{= k,} \hlkwc{ncol} \hldef{=} \hlkwd{ncol}\hldef{(X_tfidf))}
\hldef{err_history} \hlkwb{=} \hlkwd{numeric}\hldef{(}\hlnum{30}\hldef{)}

\hlkwa{for} \hldef{(iter} \hlkwa{in} \hlnum{1}\hlopt{:}\hlnum{30}\hldef{) \{}
  \hlcom{# Update H by least squares and clip negatives to 0}
  \hldef{H} \hlkwb{=} \hlkwd{pmax}\hldef{(}\hlnum{0}\hldef{,} \hlkwd{solve}\hldef{(}\hlkwd{t}\hldef{(W)} \hlopt{%*%} \hldef{W} \hlopt{+} \hlnum{1e-8} \hlopt{*} \hlkwd{diag}\hldef{(k),} \hlkwd{t}\hldef{(W)} \hlopt{%*%} \hldef{X_tfidf))}
  \hlcom{# Update W by least squares and clip negatives to 0}
  \hldef{W} \hlkwb{=} \hlkwd{pmax}\hldef{(}\hlnum{0}\hldef{,} \hlkwd{t}\hldef{(}\hlkwd{solve}\hldef{(H} \hlopt{%*%} \hlkwd{t}\hldef{(H)} \hlopt{+} \hlnum{1e-8} \hlopt{*} \hlkwd{diag}\hldef{(k), H} \hlopt{%*%} \hlkwd{t}\hldef{(X_tfidf))))}
  \hlcom{# Normalise columns of W and rescale rows of H accordingly}
  \hldef{w_norms} \hlkwb{=} \hlkwd{sqrt}\hldef{(}\hlkwd{colSums}\hldef{(W}\hlopt{^}\hlnum{2}\hldef{))}
  \hldef{W} \hlkwb{=} \hldef{W} \hlopt{%*%} \hlkwd{diag}\hldef{(}\hlnum{1} \hlopt{/} \hldef{w_norms)}
  \hldef{H} \hlkwb{=} \hlkwd{diag}\hldef{(w_norms)} \hlopt{%*%} \hldef{H}
  \hldef{err_history[iter]} \hlkwb{=} \hlkwd{sum}\hldef{((X_tfidf} \hlopt{-} \hldef{W} \hlopt{%*%} \hldef{H)}\hlopt{^}\hlnum{2}\hldef{)}
\hldef{\}}
\end{alltt}


{\ttfamily\noindent\bfseries\color{errorcolor}{\#\# Error in `h()`:\\\#\# ! error in evaluating the argument 'x' in selecting a method for function 't': error in evaluating the argument 'a' in selecting a method for function 'solve': non-conformable arrays}}\begin{alltt}
\hlkwd{rownames}\hldef{(W)} \hlkwb{=} \hldef{terms;} \hlkwd{colnames}\hldef{(W)} \hlkwb{=} \hlkwd{c}\hldef{(}\hlsng{"Topic_1"}\hldef{,} \hlsng{"Topic_2"}\hldef{)}
\hlkwd{rownames}\hldef{(H)} \hlkwb{=} \hlkwd{c}\hldef{(}\hlsng{"Topic_1"}\hldef{,} \hlsng{"Topic_2"}\hldef{);} \hlkwd{colnames}\hldef{(H)} \hlkwb{=} \hlkwd{paste0}\hldef{(}\hlsng{"d"}\hldef{,} \hlnum{1}\hlopt{:}\hlnum{6}\hldef{)}
\end{alltt}


{\ttfamily\noindent\bfseries\color{errorcolor}{\#\# Error in `rownames<-`:\\\#\# ! attempt to set 'rownames' on an object with no dimensions}}\end{kframe}
\end{knitrout}
\end{frame}


\begin{frame}[fragile]{R demonstration -- Inspecting the learned $W$ and $H$ factors}
\begin{knitrout}
\definecolor{shadecolor}{rgb}{0.969, 0.969, 0.969}\color{fgcolor}\begin{kframe}
\begin{alltt}
\hlkwd{round}\hldef{(W,} \hlnum{3}\hldef{)}
\end{alltt}
\begin{verbatim}
##        Topic_1 Topic_2
## matrix   0.918   0.818
## vector   0.420   0.933
## eigen    0.361   0.845
## rank     0.661   0.414
## graph    0.333   0.831
## vertex   0.697   0.877
## edge     0.269   0.755
## data     0.803   0.867
\end{verbatim}
\begin{alltt}
\hlkwd{round}\hldef{(H,} \hlnum{2}\hldef{)}
\end{alltt}
\begin{verbatim}
##  [1] 1.06 0.35 1.78 0.00 0.50 0.60 0.00 1.93 0.00 2.21 0.00 1.82
\end{verbatim}
\end{kframe}
\end{knitrout}
\end{frame}




\maxFrameImage{FIGS/U11-plot-nmf-results-1.pdf}


\begin{frame}{Summary of Unit U11}
\begin{itemize}
\item \textbf{Tokenisation and bag-of-words}: A vocabulary $V = \{w_1, \ldots, w_m\}$ maps each token to $\be_i \in \IR^m$ and each document $d_j$ to a sparse count vector $\bx_j = \sum_{i=1}^m f_{ij} \be_i \in \IRplus^m$ with $\|\bx_j\|_0 \ll m$
\vfill
\item \textbf{Term-document matrices and TF-IDF scaling}: Stacking document vectors as columns forms $X \in \M_{m,n}(\IRplus)$ where TF-IDF with length normalisation is the two-sided diagonal matrix product
\[
\tilde{X} = D_{\text{IDF}} X D_{\text{len}}^{-1}
\]
\vfill
\item \textbf{Cosine similarity in high dimensions}: Semantic similarity is measured by the scale-invariant angle $\mathsf{sim}(\bx, \by) = \frac{\bx \bullet \by}{\|\bx\|_2 \|\by\|_2} \in [0, 1]$ on $\IRplus^m$ whose pairwise values form the Gram matrix $S = \hat{X}^T \hat{X}$
\vfill
\item \textbf{Non-negative Matrix Factorisation (NMF)}: Approximating $X \approx W H$ with $W \in \M_{m,k}(\IRplus)$ and $H \in \M_{k,n}(\IRplus)$ avoids the subtractive cancellation of SVD, yielding additive topic columns $\bw_r$ and document encodings $\bh_j$ computed via projected \textbf{Alternating Least Squares (ALS)}
\end{itemize}
\end{frame}




\end{document}
